% Options for packages loaded elsewhere
\PassOptionsToPackage{unicode}{hyperref}
\PassOptionsToPackage{hyphens}{url}
\documentclass[
  canadien,
  11pt,
  a4paper,
]{report}
\usepackage{xcolor}
\usepackage[margin=22mm]{geometry}
\usepackage{amsmath,amssymb}
\setcounter{secnumdepth}{-\maxdimen} % remove section numbering
\usepackage{iftex}
\ifPDFTeX
  \usepackage[T1]{fontenc}
  \usepackage[utf8]{inputenc}
  \usepackage{textcomp} % provide euro and other symbols
\else % if luatex or xetex
  \usepackage{unicode-math} % this also loads fontspec
  \defaultfontfeatures{Scale=MatchLowercase}
  \defaultfontfeatures[\rmfamily]{Ligatures=TeX,Scale=1}
\fi
\usepackage{lmodern}
\ifPDFTeX\else
  % xetex/luatex font selection
  \setmainfont[]{Arial}
  \setsansfont[]{Arial}
  \setmonofont[]{Consolas}
\fi
% Use upquote if available, for straight quotes in verbatim environments
\IfFileExists{upquote.sty}{\usepackage{upquote}}{}
\IfFileExists{microtype.sty}{% use microtype if available
  \usepackage[]{microtype}
  \UseMicrotypeSet[protrusion]{basicmath} % disable protrusion for tt fonts
}{}
\makeatletter
\@ifundefined{KOMAClassName}{% if non-KOMA class
  \IfFileExists{parskip.sty}{%
    \usepackage{parskip}
  }{% else
    \setlength{\parindent}{0pt}
    \setlength{\parskip}{6pt plus 2pt minus 1pt}}
}{% if KOMA class
  \KOMAoptions{parskip=half}}
\makeatother
\usepackage{color}
\usepackage{fancyvrb}
\newcommand{\VerbBar}{|}
\newcommand{\VERB}{\Verb[commandchars=\\\{\}]}
\DefineVerbatimEnvironment{Highlighting}{Verbatim}{commandchars=\\\{\}}
% Add ',fontsize=\small' for more characters per line
\newenvironment{Shaded}{}{}
\newcommand{\AlertTok}[1]{\textcolor[rgb]{1.00,0.00,0.00}{\textbf{#1}}}
\newcommand{\AnnotationTok}[1]{\textcolor[rgb]{0.38,0.63,0.69}{\textbf{\textit{#1}}}}
\newcommand{\AttributeTok}[1]{\textcolor[rgb]{0.49,0.56,0.16}{#1}}
\newcommand{\BaseNTok}[1]{\textcolor[rgb]{0.25,0.63,0.44}{#1}}
\newcommand{\BuiltInTok}[1]{\textcolor[rgb]{0.00,0.50,0.00}{#1}}
\newcommand{\CharTok}[1]{\textcolor[rgb]{0.25,0.44,0.63}{#1}}
\newcommand{\CommentTok}[1]{\textcolor[rgb]{0.38,0.63,0.69}{\textit{#1}}}
\newcommand{\CommentVarTok}[1]{\textcolor[rgb]{0.38,0.63,0.69}{\textbf{\textit{#1}}}}
\newcommand{\ConstantTok}[1]{\textcolor[rgb]{0.53,0.00,0.00}{#1}}
\newcommand{\ControlFlowTok}[1]{\textcolor[rgb]{0.00,0.44,0.13}{\textbf{#1}}}
\newcommand{\DataTypeTok}[1]{\textcolor[rgb]{0.56,0.13,0.00}{#1}}
\newcommand{\DecValTok}[1]{\textcolor[rgb]{0.25,0.63,0.44}{#1}}
\newcommand{\DocumentationTok}[1]{\textcolor[rgb]{0.73,0.13,0.13}{\textit{#1}}}
\newcommand{\ErrorTok}[1]{\textcolor[rgb]{1.00,0.00,0.00}{\textbf{#1}}}
\newcommand{\ExtensionTok}[1]{#1}
\newcommand{\FloatTok}[1]{\textcolor[rgb]{0.25,0.63,0.44}{#1}}
\newcommand{\FunctionTok}[1]{\textcolor[rgb]{0.02,0.16,0.49}{#1}}
\newcommand{\ImportTok}[1]{\textcolor[rgb]{0.00,0.50,0.00}{\textbf{#1}}}
\newcommand{\InformationTok}[1]{\textcolor[rgb]{0.38,0.63,0.69}{\textbf{\textit{#1}}}}
\newcommand{\KeywordTok}[1]{\textcolor[rgb]{0.00,0.44,0.13}{\textbf{#1}}}
\newcommand{\NormalTok}[1]{#1}
\newcommand{\OperatorTok}[1]{\textcolor[rgb]{0.40,0.40,0.40}{#1}}
\newcommand{\OtherTok}[1]{\textcolor[rgb]{0.00,0.44,0.13}{#1}}
\newcommand{\PreprocessorTok}[1]{\textcolor[rgb]{0.74,0.48,0.00}{#1}}
\newcommand{\RegionMarkerTok}[1]{#1}
\newcommand{\SpecialCharTok}[1]{\textcolor[rgb]{0.25,0.44,0.63}{#1}}
\newcommand{\SpecialStringTok}[1]{\textcolor[rgb]{0.73,0.40,0.53}{#1}}
\newcommand{\StringTok}[1]{\textcolor[rgb]{0.25,0.44,0.63}{#1}}
\newcommand{\VariableTok}[1]{\textcolor[rgb]{0.10,0.09,0.49}{#1}}
\newcommand{\VerbatimStringTok}[1]{\textcolor[rgb]{0.25,0.44,0.63}{#1}}
\newcommand{\WarningTok}[1]{\textcolor[rgb]{0.38,0.63,0.69}{\textbf{\textit{#1}}}}
\usepackage{longtable,booktabs,array}
\newcounter{none} % for unnumbered tables
\usepackage{calc} % for calculating minipage widths
% Correct order of tables after \paragraph or \subparagraph
\usepackage{etoolbox}
\makeatletter
\patchcmd\longtable{\par}{\if@noskipsec\mbox{}\fi\par}{}{}
\makeatother
% Allow footnotes in longtable head/foot
\IfFileExists{footnotehyper.sty}{\usepackage{footnotehyper}}{\usepackage{footnote}}
\makesavenoteenv{longtable}
\ifLuaTeX
\usepackage[bidi=basic,shorthands=off]{babel}
\else
\usepackage[bidi=default,shorthands=off]{babel}
\fi
\ifPDFTeX
\else
\babelfont{rm}[]{Arial}
\fi
\ifLuaTeX
  \usepackage{selnolig} % disable illegal ligatures
\fi
\setlength{\emergencystretch}{3em} % prevent overfull lines
\providecommand{\tightlist}{%
  \setlength{\itemsep}{0pt}\setlength{\parskip}{0pt}}
\usepackage{fvextra}
\DefineVerbatimEnvironment{Highlighting}{Verbatim}{breaklines,breakanywhere,commandchars=\\\{\}}
\fvset{fontsize=\small}
\setlength{\emergencystretch}{3em}
\definecolor{orbitcyan}{HTML}{006C82}
\AtBeginDocument{\renewcommand{\contentsname}{Table des matières}\hypersetup{colorlinks=true,linkcolor=orbitcyan,urlcolor=orbitcyan}}
\usepackage{bookmark}
\IfFileExists{xurl.sty}{\usepackage{xurl}}{} % add URL line breaks if available
\urlstyle{same}
\hypersetup{
  pdftitle={Orbit Formation},
  pdfauthor={Jean-Sébastien Beaulieu},
  pdflang={fr-CA},
  hidelinks,
  pdfcreator={LaTeX via pandoc}}

\title{Orbit Formation}
\usepackage{etoolbox}
\makeatletter
\providecommand{\subtitle}[1]{% add subtitle to \maketitle
  \apptocmd{\@title}{\par {\large #1 \par}}{}{}
}
\makeatother
\subtitle{Construire, vérifier et enseigner avec son assistant
personnel}
\author{Jean-Sébastien Beaulieu}
\date{}

\begin{document}
\maketitle

{
\setcounter{tocdepth}{2}
\tableofcontents
}
\chapter{Orbit Formation}\label{orbit-formation}

\section{Construire, vérifier et enseigner avec son assistant
personnel}\label{construire-vuxe9rifier-et-enseigner-avec-son-assistant-personnel}

Jean-Sébastien Beaulieu · Édition pédagogique 1.0 · Français

Ce manuel accompagne une formation de \textbf{40 heures}. Sa première
partie aide à comprendre et pratiquer les mécanismes. Sa seconde partie
aide l'enseignant à les faire observer, expliquer et transférer. Les
fichiers du cours, les interventions de l'assistant et les modifications
de l'élève gardent une attribution distincte.

\subsection{Note d'édition et de
responsabilité}\label{note-duxe9dition-et-de-responsabilituxe9}

Le code et les documents possèdent des versions différentes :
\texttt{orbit-course-1.0.0} identifie le cours ; la version d'un landing
ou d'un dossier désigne son propre objet. Ce document est une source
préparée pour revue. La disponibilité d'un PDF, d'une page Canva ou d'un
outil n'est pas un résultat d'apprentissage. Les conditions techniques
sont réunies dans la section de qualification et dans le reçu maintenu.

Les droits des ressources citées restent ceux de leurs auteurs et de
leurs licences. Cette édition identifie l'intention pédagogique et ses
sources ; elle n'attribue aucune nouvelle licence aux documents
externes. Le partage des travaux d'un élève suit son choix et les droits
de son environnement.

\subsection{Collaboration --- déclaration à examiner lors de la revue
d'auteur}\label{collaboration-duxe9claration-uxe0-examiner-lors-de-la-revue-dauteur}

J'ai utilisé Codex à son plein potentiel comme partenaire de recherche
--- pour la cartographie du code, la comparaison des sources,
l'organisation des preuves, les contrôles de cohérence, le contrôle
éditorial et la préparation des livrables. J'ai formulé l'intention,
défini le périmètre, interprété les résultats, arbitré les conclusions
et conservé chaque décision publique. Cette collaboration élargit ma
capacité d'investigation ; le jugement, la responsabilité, la qualité
d'auteur et la signature finale restent sous mon autorité.

\subsection{À qui s'adresse le manuel
?}\label{uxe0-qui-sadresse-le-manuel}

À un élève qui apprend avec son assistant personnel, et à l'enseignant
qui accompagne ses essais. Les notions sont introduites depuis une
action visible et un fichier réel. Un exemple complet peut être demandé
à l'assistant : l'élève vérifie ensuite ce qu'il conserve, l'explique ou
l'applique à une autre situation.

\subsection{Comment l'utiliser}\label{comment-lutiliser}

Lire la fiche du module, retrouver ses fichiers et poser une prédiction.
Exécuter l'activité Colab avec une modification contrôlée. Choisir un
rendu à remettre. Pendant la séance, expliquer une observation et
essayer une situation nouvelle. La seconde partie du manuel fournit les
scripts et les pistes de revue à l'enseignant ; les documents web élève
se distribuent au rythme du cours.

Les mots \textbf{fourni, modifié, déclaré exécuté, vérifié techniquement
et examiné humainement} désignent des étapes différentes. Une aide peut
être expliquée franchement. Une question ouverte et une limite donnent à
l'enseignant un point de départ concret.

\section{Table des matières}\label{table-des-matiuxe8res}

\begin{enumerate}
\def\labelenumi{\arabic{enumi}.}
\tightlist
\item
  Le parcours de 40 heures et les destinations du travail.
\item
  Partie I --- apprendre : huit modules et projet personnel.
\item
  Partie II --- enseigner : huit fiches de séance et projet.
\item
  Annexes : capacités, remise, code, glossaire, symboles et sources.
\end{enumerate}

La table des matières paginée sera construite lors de l'export DOCX/PDF.

\chapter{Le parcours de 40 heures}\label{le-parcours-de-40-heures}

{\def\LTcaptype{none} % do not increment counter
\begin{longtable}[]{@{}lrrr@{}}
\toprule\noalign{}
Partie & Avec l'enseignant & Solo avec l'assistant personnel & Total \\
\midrule\noalign{}
\endhead
\bottomrule\noalign{}
\endlastfoot
Huit modules & 8 × 1 h & 8 × 3 h & 32 h \\
Projet personnel & 2 h & 6 h & 8 h \\
\textbf{Formation complète} & \textbf{10 h} & \textbf{30 h} & \textbf{40
h} \\
\end{longtable}
}

Chaque période solo comprend \textbf{120 minutes de lecture et
d'apprentissage guidé, 30 minutes d'activité Colab, puis 30 minutes de
bilan}. Le webinaire ajoute une heure déjà comptée dans le module : 10
minutes de revue, 15 de clarification, 25 de construction et 10 de
transfert.

Le projet suit \textbf{1 h de cadrage accompagné → 6 h solo → 1 h de
clôture accompagnée}. Ses heures sont distinctes des modules. Dans le
calendrier de seize jours, les six heures solo se distribuent comme 1 +
2 + 2 + 1 pendant les quatre derniers modules. La clôture suit les six
heures solo.

\section{Le rôle de chacun}\label{le-ruxf4le-de-chacun}

{\def\LTcaptype{none} % do not increment counter
\begin{longtable}[]{@{}
  >{\raggedright\arraybackslash}p{(\linewidth - 2\tabcolsep) * \real{0.5000}}
  >{\raggedright\arraybackslash}p{(\linewidth - 2\tabcolsep) * \real{0.5000}}@{}}
\toprule\noalign{}
\begin{minipage}[b]{\linewidth}\raggedright
Personne ou système
\end{minipage} & \begin{minipage}[b]{\linewidth}\raggedright
Travail et décision
\end{minipage} \\
\midrule\noalign{}
\endhead
\bottomrule\noalign{}
\endlastfoot
Élève & Choisir son intention, ses essais, le niveau d'aide et les
éléments à remettre ; expliquer sa décision \\
Assistant personnel & Expliquer, suggérer et comparer ; conserver les
conditions, sources et questions ouvertes \\
Enseignant & Définir la progression, examiner le travail, choisir un
exemple utile et enregistrer une revue réellement effectuée \\
Orbit & Fournir ressources, états, calculs déterministes, partage
sélectionné et exports \\
\end{longtable}
}

Une réponse complète de l'assistant peut être utile. La question
pédagogique devient : comment cette réponse est-elle vérifiée, expliquée
et réutilisée ? Les critères portent sur des gestes observables et des
raisonnements reliés au code.

\section{Trois destinations à
comprendre}\label{trois-destinations-uxe0-comprendre}

{\def\LTcaptype{none} % do not increment counter
\begin{longtable}[]{@{}
  >{\raggedright\arraybackslash}p{(\linewidth - 2\tabcolsep) * \real{0.5000}}
  >{\raggedright\arraybackslash}p{(\linewidth - 2\tabcolsep) * \real{0.5000}}@{}}
\toprule\noalign{}
\begin{minipage}[b]{\linewidth}\raggedright
Destination
\end{minipage} & \begin{minipage}[b]{\linewidth}\raggedright
Contenu et choix
\end{minipage} \\
\midrule\noalign{}
\endhead
\bottomrule\noalign{}
\endlastfoot
Orbit en mémoire ou dans le profil navigateur & Session et fichiers
sélectionnés ; la sauvegarde locale est un choix explicite, avec export
avant fermeture si elle reste indisponible \\
Colab / Drive personnel & Copie du carnet, code, texte et sorties ; il
s'agit du cloud Google et le partage peut inclure ces éléments \\
Sanity & Contenu accepté comme publiable, après examen du contenu et de
la destination dans le plugin Studio \\
\end{longtable}
}

Le journal privé et la conversation complète ne sont pas des pièces
obligatoires de remise. Partager avec l'assistant, déposer une
proposition, choisir un moteur, remettre à l'enseignant et publier
restent des décisions distinctes. Un stockage dans le profil navigateur
n'est pas présenté comme un coffre chiffré.

\section{Un notebook gratuit et une reprise
explicite}\label{un-notebook-gratuit-et-une-reprise-explicite}

Les exercices obligatoires utilisent un runtime CPU de Colab, aucun GPU,
aucune API payante et aucune fonctionnalité d'IA intégrée obligatoire.
Le Module 3 charge Three.js dans le navigateur depuis un CDN avec une
alternative HTML. Les ressources Colab gratuites restent variables :
conserver les exports et savoir reprendre fait partie du parcours.

Le code dans \texttt{student\_\allowbreak{}files} passe à l'aperçu puis au ZIP. Le
résultat JSON conserve prédiction, paramètres, observations choisies,
explication, aide reçue, limites et question ouverte. Le ZIP contient le
code choisi, une copie rejouable du notebook, les instructions et un
manifeste SHA-256. Son import conserve le statut
\texttt{external-declared}.

La fiche de Module 2 distingue aperçu et compilation Astro. Celle du
Module 8 distingue trace préparée et appel natif. Les vérifications sur
Colab, Orbit, un projet assemblé et un second Studio conservent chacune
leur environnement.

\chapter{Partie I --- comprendre, expérimenter et
construire}\label{partie-i-comprendre-expuxe9rimenter-et-construire}

Le parcours commun est \textbf{observer → questionner → expliquer →
prédire → expérimenter → conserver une preuve → transférer}. Ouvrir le
fichier indiqué avant de modifier sa valeur. L'enseignant pourra
examiner un texte, un schéma ou du code annoté ; une forme unique de
réponse n'est pas imposée.

\chapter{Module 1 --- Geste, état et
rendu}\label{module-1-geste-uxe9tat-et-rendu}

\textbf{Objectif : expliquer quelle donnée un événement modifie et
comment le rendu en dépend.} À la fin de l'activité, vous devez pouvoir
retrouver la position, le pointeur actif et la phase du geste dans le
code, puis expliquer un relâchement et une interruption. La fiche
fournit une méthode ; votre compréhension sera examinée pendant le
webinaire.

Ce module occupe les jours 1--2 : \textbf{3 h solo, puis 1 h avec
l'enseignant}. Le notebook et le modèle sont fournis. Vos modifications
et votre explication restent attribuées séparément à vous et à votre
assistant.

\section{Avant de commencer}\label{avant-de-commencer}

Ouvrez le \href{../../notebooks/module-1.ipynb}{notebook élève} dans
Colab, puis conservez une copie personnelle. Un compte Google gratuit et
un runtime CPU suffisent. L'assistant personnel peut fournir un exemple
complet ; vous choisissez ce que vous essayez et remettez.

Le code d'étude est
\href{../../frontend/module-1/interaction-state.js}{interaction-state.js}.
La copie de correction est conservée dans l'espace enseignant. La
mission du catalogue est définie dans
\href{../../../../../packages/learning/src/catalog.ts}{catalog.ts}.

\section{Deux heures d'apprentissage
guidé}\label{deux-heures-dapprentissage-guiduxe9}

{\def\LTcaptype{none} % do not increment counter
\begin{longtable}[]{@{}
  >{\raggedright\arraybackslash}p{(\linewidth - 4\tabcolsep) * \real{0.3333}}
  >{\raggedright\arraybackslash}p{(\linewidth - 4\tabcolsep) * \real{0.3333}}
  >{\raggedright\arraybackslash}p{(\linewidth - 4\tabcolsep) * \real{0.3333}}@{}}
\toprule\noalign{}
\begin{minipage}[b]{\linewidth}\raggedright
Temps
\end{minipage} & \begin{minipage}[b]{\linewidth}\raggedright
Votre travail avec l'assistant
\end{minipage} & \begin{minipage}[b]{\linewidth}\raggedright
Trace sélectionnée
\end{minipage} \\
\midrule\noalign{}
\endhead
\bottomrule\noalign{}
\endlastfoot
0--30 min & Lire le modèle. Repérer \texttt{state.position},
\texttt{phase} et \texttt{pointer\allowbreak{}Id}. Dessiner « événement → donnée →
rendu ». & Un petit schéma et une question \\
30--60 min & Lire \texttt{pointerdown} et \texttt{pointermove}.
Expliquer pourquoi un survol seul ne déplace pas l'élément. & Trois
phrases reliées aux lignes concernées \\
60--90 min & Comparer \texttt{pointerup}, \texttt{pointercancel} et
\texttt{lostpointercapture}. Lire les ressources officielles. & Une
différence entre relâchement et annulation \\
90--120 min & Reprendre le parcours au clavier. Préparer une prédiction
avec deux pressions vers la droite. & Valeur initiale, valeur testée,
conditions \\
\end{longtable}
}

Demande possible à l'assistant : « Explique-moi \texttt{pointer\allowbreak{}Id} à
partir de ce fichier. Donne un scénario normal et un scénario
interrompu, puis demande-moi quelle donnée doit être restaurée. » Vous
pouvez demander la solution ; votre propre trace décrit ensuite ce que
vous avez vérifié.

Dans ce modèle, \texttt{x} et \texttt{y} sont normalisés entre 0 et 1.
\texttt{x\ =\ 0.5} signifie le milieu de la zone, pas 0,5 pixel. L'état
est copié avant transmission au rendu, afin de garder une séparation
explicite entre les données et leur présentation.

\begin{Shaded}
\begin{Highlighting}[]
\ImportTok{export} \KeywordTok{const}\NormalTok{ interactionOptions }\OperatorTok{=}\NormalTok{ \{ }\DataTypeTok{keyboardStep}\OperatorTok{:} \FloatTok{0.04}\NormalTok{ \}}\OperatorTok{;}
\end{Highlighting}
\end{Shaded}

Cette valeur règle le pas clavier. Elle ne change ni la largeur de la
zone ni la vitesse de la souris. Sous capture du pointeur, sortir de la
zone ne signifie pas automatiquement annuler le geste.

\section{Activité Colab --- 30
minutes}\label{activituxe9-colab-30-minutes}

\begin{enumerate}
\def\labelenumi{\arabic{enumi}.}
\tightlist
\item
  \textbf{Prédire --- 5 min.} À partir de \texttt{x\ =\ 0.5}, annoncer
  la position attendue après deux flèches droites avec un pas de 0,04,
  puis 0,08. Ajouter une hypothèse sur le pointeur actif après
  annulation. Écrire la prédiction avant l'aperçu.
\item
  \textbf{Modifier et observer --- 15 min.} Dans
  \texttt{student\_\allowbreak{}files}, changer uniquement
  \texttt{keyboard\allowbreak{}Step:\ 0.04} en \texttt{keyboard\allowbreak{}Step:\ 0.08}. Exécuter
  l'aperçu depuis cette même variable. Comparer deux pressions ; essayer
  un geste puis son relâchement. Repérer dans le code la transition
  d'annulation. Une annulation native dépend du navigateur et du
  périphérique ; ne la déclarez pas observée si vous avez seulement lu
  le code ou pressé Échap.
\item
  \textbf{Expliquer et exporter --- 10 min.} Noter le résultat observé,
  la variable conservée et une limite. Remplir l'aide reçue et
  l'explication, puis exporter le ZIP. Le ZIP contient le code
  effectivement présent dans \texttt{student\_\allowbreak{}files}, pas un corrigé
  substitué.
\end{enumerate}

Piste de vérification, à consulter après votre prédiction : deux
pressions ajoutent 0,08 avec le premier pas et 0,16 avec le second, tant
que la borne droite n'est pas atteinte. Le code place \texttt{pointer\allowbreak{}Id}
à \texttt{null} à la fin d'un geste. Ce calcul attendu ne remplace pas
votre essai.

\section{Bilan solo --- 30 minutes}\label{bilan-solo-30-minutes}

Consacrer 10 minutes au choix des traces à remettre, 10 minutes à une
explication personnelle et 10 minutes à la question pour le webinaire.
Conserver le ZIP, le scénario utilisé et les éventuels écarts. Une
capture seule montre une image ; l'explication doit nommer l'événement
et l'état.

Votre carnet personnel et ses sorties sont dans Colab/Drive. Vous
choisissez le rendu transmis à l'enseignant. Aucune conversation privée
complète n'est requise.

\section{Transfert et critères
observables}\label{transfert-et-crituxe8res-observables}

Créer une carte déplaçable qui dispose aussi d'une commande clavier.
Vous devez pouvoir : nommer la donnée modifiée ; distinguer relâchement,
sortie de zone et annulation ; montrer la remise à zéro du pointeur
actif ; sélectionner le même élément sans souris. Un scénario vérifié et
une reformulation valent davantage qu'une définition copiée.

Les \href{../../../../../tests/learning-notebooks-bricks.test.ts}{tests
logiciels} examinent des comportements du modèle. Ils ne certifient ni
un geste tactile réel ni votre compréhension.

\textbf{Ressources primaires et limites :}
\href{https://developer.mozilla.org/en-US/docs/Web/API/Element/setPointerCapture}{Capture
du pointeur --- MDN} ·
\href{https://developer.mozilla.org/en-US/docs/Web/API/Element/pointercancel_event}{Événement
pointercancel --- MDN} · \href{../../PROJECT.md}{Projet et règles de
remise}

\chapter{Module 2 --- Organisation
Astro}\label{module-2-organisation-astro}

\textbf{Objectif : séparer contenu, présentation et comportement.} Vous
apprendrez à retrouver le fichier responsable d'une modification et à
réutiliser une carte avec un autre contenu. Le module occupe les jours
3--4 : \textbf{3 h solo, puis 1 h avec l'enseignant}.

Le prérequis est d'avoir rencontré le parcours événement → état → rendu
du \href{../../modules/module-1.md}{Module 1}. Votre assistant peut
expliquer directement Astro ; votre travail consiste ensuite à vérifier
cette explication dans les fichiers et à préparer une modification
précise.

\section{Fichiers à examiner}\label{fichiers-uxe0-examiner}

\begin{itemize}
\tightlist
\item
  \href{../../frontend/module-2/ExplorationCard.astro}{ExplorationCard.astro}
  fournit le titre, la section, la liste et le paragraphe de
  description.
\item
  \href{../../frontend/module-2/exploration-card.js}{exploration-card.js}
  crée les boutons, associe un identifiant à chaque élément et actualise
  la description.
\item
  \href{../../notebooks/module-2.ipynb}{Notebook élève} ; la copie de
  correction conservée par l'enseignant.
\end{itemize}

La carte Astro déclare deux propriétés : \texttt{id} et \texttt{title}.
Sa liste HTML est vide au départ ; le contrôleur fourni y ajoute les
éléments. Cette séparation permet de changer les données sans recopier
les événements.

\begin{Shaded}
\begin{Highlighting}[]
\ImportTok{export} \KeywordTok{const}\NormalTok{ cardOptions }\OperatorTok{=}\NormalTok{ \{ }\DataTypeTok{descriptionPrefix}\OperatorTok{:} \StringTok{\textquotesingle{}Mon observation : \textquotesingle{}}\NormalTok{ \}}\OperatorTok{;}
\end{Highlighting}
\end{Shaded}

Le préfixe appartient à la description affichée par le contrôleur.
Modifier ce texte ne doit pas ajouter un bouton. Dans le fichier Astro,
\texttt{aria-label} reprend le titre de la section : un libellé reste
nécessaire même si la carte change d'apparence.

\section{Deux heures d'apprentissage
guidé}\label{deux-heures-dapprentissage-guiduxe9-1}

{\def\LTcaptype{none} % do not increment counter
\begin{longtable}[]{@{}
  >{\raggedright\arraybackslash}p{(\linewidth - 4\tabcolsep) * \real{0.3333}}
  >{\raggedright\arraybackslash}p{(\linewidth - 4\tabcolsep) * \real{0.3333}}
  >{\raggedright\arraybackslash}p{(\linewidth - 4\tabcolsep) * \real{0.3333}}@{}}
\toprule\noalign{}
\begin{minipage}[b]{\linewidth}\raggedright
Temps
\end{minipage} & \begin{minipage}[b]{\linewidth}\raggedright
Travail avec l'assistant
\end{minipage} & \begin{minipage}[b]{\linewidth}\raggedright
Trace sélectionnée
\end{minipage} \\
\midrule\noalign{}
\endhead
\bottomrule\noalign{}
\endlastfoot
0--30 min & Lire le composant Astro. Retrouver propriétés, section,
titre, liste et description. & Arborescence courte et rôle de chaque
élément \\
30--60 min & Lire \texttt{create\allowbreak{}Cards}, \texttt{select} et
\texttt{dispose}. Suivre un identifiant depuis les données jusqu'au
bouton. & Un parcours de sélection annoté \\
60--90 min & Lire les composants Astro et les îlots. Distinguer page
construite, HTML reçu et comportement navigateur. & Explication de trois
responsabilités \\
90--120 min & Choisir un autre contenu pour une carte et préparer une
modification unique. & Prédiction et contenu personnel d'essai \\
\end{longtable}
}

Demande possible : « Voici ces deux fichiers. Explique qui fournit le
HTML, qui remplit la liste et qui réagit à la sélection. Donne un
exemple complet, puis demande-moi où je changerais le titre sans changer
l'état du geste. » Votre assistant aide à lire ; les chemins et symboles
servent à contrôler sa réponse.

\section{Activité Colab --- 30
minutes}\label{activituxe9-colab-30-minutes-1}

\begin{enumerate}
\def\labelenumi{\arabic{enumi}.}
\tightlist
\item
  \textbf{Prédire --- 5 min.} Écrire ce qui changera si
  \texttt{description\allowbreak{}Prefix} devient
  \texttt{Ma\ décision\ expliquée\ :}. Prédire aussi ce qui restera
  identique : nombre de boutons, identifiants et sélection.
\item
  \textbf{Modifier et observer --- 15 min.} Dans
  \texttt{student\_\allowbreak{}files}, changer seulement le préfixe de
  \texttt{exploration-card.js}. Exécuter l'aperçu, sélectionner deux
  éléments et comparer la description. Retrouver ensuite le titre dans
  \texttt{Exploration\allowbreak{}Card.astro}, sans attribuer la modification au
  mauvais fichier.
\item
  \textbf{Expliquer et exporter --- 10 min.} Décrire la responsabilité
  de chaque fichier ; conserver votre aide reçue, une limite et une
  question. Exporter les deux fichiers dans le même ZIP.
\end{enumerate}

L'aperçu du notebook présente une contrepartie HTML préparée. \textbf{Il
ne compile pas le composant Astro.} L'exercice permet de comparer le
comportement de la carte ; la compilation du vrai composant appartient
au parcours de projet vérifié séparément.

Piste après tentative : \texttt{create\allowbreak{}Cards} utilise
\texttt{text\allowbreak{}Content} pour les titres et descriptions, puis
\texttt{aria-pressed} pour la sélection. Rechercher \texttt{select(id)}
permet de relier le bouton aux données. Une amélioration de style se
vérifie ensuite dans le projet, avec le titre et le clavier préservés.

\section{Bilan solo --- 30 minutes}\label{bilan-solo-30-minutes-1}

Prendre 10 minutes pour choisir votre modification et les traces, 10
minutes pour expliquer la séparation des responsabilités et 10 minutes
pour préparer le webinaire. Votre rendu comprend le ZIP, un exemple de
carte pour votre sujet et une phrase qui distingue aperçu HTML et
compilation Astro.

Une suggestion de l'assistant, un code exécuté et une explication
comprise ont des statuts distincts. Vous pouvez signaler « je retrouve
le fichier, mais je ne peux pas encore expliquer ce contrôleur » : cette
limite prépare une revue utile.

\section{Transfert et critères
observables}\label{transfert-et-crituxe8res-observables-1}

Préparer une carte pour un thème personnel. L'élève retrouve le rôle de
chaque fichier, change une propriété sans dupliquer la logique et
explique la différence entre l'aperçu et la compilation. La sélection
doit garder un libellé et un état accessibles.

Le squelette et les contrôleurs fournis sont attribués au cours. Votre
contenu, vos modifications et les contributions de l'assistant sont
identifiés dans le bilan.

\textbf{Ressources et statut :}
\href{https://docs.astro.build/en/basics/astro-components/}{Composants
Astro --- documentation officielle} ·
\href{https://docs.astro.build/en/concepts/islands/}{Architecture en
îlots --- documentation officielle} ·
\href{../../../../../packages/learning/src/catalog.ts}{Mission et
critères du catalogue} · \href{../../PROJECT.md}{Règles du projet et de
la remise}

\chapter{Module 3 --- Scène Three.js et
coordonnées}\label{module-3-scuxe8ne-three.js-et-coordonnuxe9es}

\textbf{Objectif : relier objet, caméra, coordonnées et sélection.} Vous
allez observer comment la caméra modifie une projection, puis retrouver
le même élément dans une représentation HTML. Le module occupe les jours
5--6 : \textbf{3 h solo, puis 1 h avec l'enseignant}.

Le prérequis est de distinguer structure et comportement dans une carte.
La scène est une autre manière de présenter les éléments ; leur
identifiant et leur contenu restent disponibles dans la liste HTML.

\section{Code et matériel}\label{code-et-matuxe9riel}

Ouvrir le \href{../../notebooks/module-3.ipynb}{notebook élève} et
\href{../../frontend/module-3/scene-controller.js}{scene-controller.js}.
La copie de correction conservée par l'enseignant est séparée. Le
notebook utilise un runtime CPU ; Three.js 0.181.2 s'exécute dans le
navigateur depuis un CDN. Le refus du réseau ou l'absence de WebGL
conserve le parcours HTML, sans valider une scène 3D.

\begin{Shaded}
\begin{Highlighting}[]
\ImportTok{export} \KeywordTok{const}\NormalTok{ sceneOptions }\OperatorTok{=}\NormalTok{ \{ }\DataTypeTok{cameraDistance}\OperatorTok{:} \DecValTok{5}\NormalTok{ \}}\OperatorTok{;}
\end{Highlighting}
\end{Shaded}

Le contrôleur reçoit \texttt{THREE} depuis son hôte. Il prépare scène,
caméra, objets et lumière. Il ne crée pas sa propre boucle de rendu :
l'hôte décide quand appeler \texttt{render()}. Le même \texttt{item.id}
relie sphère, sélection et fiche HTML.

Les données de position de l'activité sont normalisées. Le contrôleur
les transforme vers la scène ; la caméra projette ensuite la scène vers
l'écran. Déplacer une caméra n'est donc pas déplacer les données de
l'objet.

\section{Deux heures d'apprentissage
guidé}\label{deux-heures-dapprentissage-guiduxe9-2}

{\def\LTcaptype{none} % do not increment counter
\begin{longtable}[]{@{}
  >{\raggedright\arraybackslash}p{(\linewidth - 4\tabcolsep) * \real{0.3333}}
  >{\raggedright\arraybackslash}p{(\linewidth - 4\tabcolsep) * \real{0.3333}}
  >{\raggedright\arraybackslash}p{(\linewidth - 4\tabcolsep) * \real{0.3333}}@{}}
\toprule\noalign{}
\begin{minipage}[b]{\linewidth}\raggedright
Temps
\end{minipage} & \begin{minipage}[b]{\linewidth}\raggedright
Travail avec l'assistant
\end{minipage} & \begin{minipage}[b]{\linewidth}\raggedright
Trace sélectionnée
\end{minipage} \\
\midrule\noalign{}
\endhead
\bottomrule\noalign{}
\endlastfoot
0--30 min & Lire scène, caméra, matériau et lumière. Associer chaque mot
à un objet du code. & Schéma de la présentation \\
30--60 min & Suivre un identifiant de la liste HTML à une sphère.
Repérer le raycasting. & Chemin « clic → identifiant → fiche » \\
60--90 min & Lire la projection et le redimensionnement. Distinguer
coordonnées normalisées, scène et écran. & Une explication avec trois
coordonnées différentes \\
90--120 min & Préparer la comparaison caméra 5/7 avec positions
conservées. Examiner le parcours sans WebGL. & Prédiction et conditions
de l'essai \\
\end{longtable}
}

Demande possible : « Donne-moi une analogie de la caméra, puis relie-la
à \texttt{camera.position.z} et à \texttt{item.position}. Demande-moi
quelle valeur change si l'objet paraît plus petit. » L'analogie aide à
démarrer ; l'explication finale revient aux données du code.

\section{Activité Colab --- 30
minutes}\label{activituxe9-colab-30-minutes-2}

\begin{enumerate}
\def\labelenumi{\arabic{enumi}.}
\tightlist
\item
  \textbf{Prédire --- 5 min.} Avec les mêmes objets, annoncer l'effet
  attendu d'une caméra passant de 5 à 7. Préciser si votre hypothèse
  concerne la taille affichée ou la position enregistrée.
\item
  \textbf{Modifier et observer --- 15 min.} Dans
  \texttt{student\_\allowbreak{}files}, remplacer \texttt{camera\allowbreak{}Distance:\ 5} par
  \texttt{camera\allowbreak{}Distance:\ 7}. Exécuter l'aperçu. Choisir une sphère
  puis le bouton HTML correspondant ; comparer l'identifiant
  sélectionné. Observer un changement de taille de la zone si
  l'interface le permet. Si la scène ne démarre pas, consigner le refus
  et utiliser la liste, sans inventer une observation 3D.
\item
  \textbf{Expliquer et exporter --- 10 min.} Décrire ce qui a changé et
  ce qui reste dans les données. Noter navigateur, disponibilité WebGL,
  aide reçue et limite. Exporter le contrôleur réellement modifié.
\end{enumerate}

Piste après tentative : reculer la caméra modifie la projection, avec
les positions d'objet inchangées. Une sphère et un bouton HTML peuvent
désigner le même élément. Un identifiant commun ne suffit toutefois pas
à prouver que tous les gestes tactiles ou tous les redimensionnements
fonctionnent.

\section{Bilan solo --- 30 minutes}\label{bilan-solo-30-minutes-2}

Choisir les traces pendant 10 minutes, rédiger votre explication pendant
10 minutes et préparer une question pendant 10 minutes. Votre rendu
nomme l'élément choisi, les deux distances, la disponibilité de la scène
et une limite.

La position, la couleur ou la taille d'un objet \textbf{ne mesure pas la
vérité d'une affirmation}. La scène fournit ici une représentation
artistique et pédagogique ; elle ne constitue pas une simulation
scientifique validée.

\section{Transfert et critères
observables}\label{transfert-et-crituxe8res-observables-2}

Créer trois éléments consultables à la fois dans une liste et une scène.
L'élève distingue coordonnées d'écran et de scène, montre le lien d'une
sélection à sa fiche et retrouve le contenu sans WebGL. Il explique le
rôle d'un redimensionnement et les ressources à nettoyer à la fermeture.

Le contrôleur possède \texttt{dispose()} pour libérer événements,
géométries, matériaux et renderer. La présence de cette fonction dans le
code demande encore une vérification de son appel dans le parcours
complet.

\textbf{Ressources et statut :}
\href{https://threejs.org/manual/\#en/fundamentals}{Fondamentaux
Three.js --- manuel officiel} ·
\href{https://threejs.org/docs/\#api/en/core/Raycaster}{Raycaster ---
API Three.js} ·
\href{../../../../../packages/learning/src/catalog.ts}{Mission du
catalogue} · \href{../../PROJECT.md}{Projet et règles de remise}

\section{Quand deux interactions partagent le même
geste}\label{quand-deux-interactions-partagent-le-muxeame-geste}

Dans
\href{../../frontend/module-3/scene-controller.js}{scene-controller.js},
\texttt{pointer\allowbreak{}Target} indique où écouter le geste. Sa valeur par défaut
est le canvas de Three.js. Le calcul de projection utilise toujours le
rectangle du canvas : la cible d'événement et la surface dessinée ont
des responsabilités différentes.

Le Module 1 peut capturer le pointeur sur le stage. Dans cet assemblage,
le relâchement est alors envoyé au stage. Un écouteur placé seulement
sur le canvas risque de ne jamais le recevoir, même si la scène se
dessine correctement. L'hôte fourni dans
\href{../../assembly/src/student-frontend.js}{student-frontend.js}
raccorde explicitement les deux briques :

\begin{Shaded}
\begin{Highlighting}[]
\NormalTok{scene }\OperatorTok{=} \FunctionTok{createSceneController}\NormalTok{(\{}
  \DataTypeTok{container}\OperatorTok{:} \BuiltInTok{root}\OperatorTok{.}\FunctionTok{querySelector}\NormalTok{(}\StringTok{\textquotesingle{}[data{-}three]\textquotesingle{}}\NormalTok{)}\OperatorTok{,}
\NormalTok{  THREE}\OperatorTok{,}\NormalTok{ items}\OperatorTok{,} \DataTypeTok{onSelect}\OperatorTok{:}\NormalTok{ select}\OperatorTok{,} \DataTypeTok{pointerTarget}\OperatorTok{:}\NormalTok{ stage}
\NormalTok{\})}\OperatorTok{;}
\end{Highlighting}
\end{Shaded}

Ce code est fourni par le cours. Il ne remplace pas le fichier exporté
de l'élève. Le début du geste doit provenir du canvas ou de son chemin
d'événements ; un déplacement supérieur à six pixels est considéré comme
un glisser et ne sélectionne pas une sphère. Une annulation efface le
geste en attente. Le nettoyage retire les écouteurs, l'observateur de
taille et les ressources de rendu.

\textbf{Petit transfert :} dessiner deux rectangles, le stage et son
canvas. Indiquer où le pointeur commence, quelle surface le capture, où
arrive son relâchement et quel rectangle transforme ses coordonnées.
Puis expliquer ce que l'on doit conserver lorsqu'une caméra recule. Les
essais réellement observés, les incidents et le retest de l'assemblage
gardent leurs reçus dans la qualification ; cet exemple de code seul
n'en démontre pas le comportement.

\chapter{Module 4 --- Vitesse et
inertie}\label{module-4-vitesse-et-inertie}

\textbf{Objectif : comprendre le mouvement après relâchement.} Vous
allez définir position, vitesse et durée, puis comparer deux
amortissements. Le module occupe les jours 7--8 : \textbf{3 h solo, puis
1 h avec l'enseignant}.

Le geste, ses phases et les coordonnées ont été rencontrés. La nouvelle
question est : quelle donnée permet à un élément de continuer son
mouvement lorsque le pointeur n'est plus actif ?

\section{Du geste au mouvement}\label{du-geste-au-mouvement}

Le modèle \href{../../frontend/module-4/inertia.js}{inertia.js} expose
\texttt{grab}, \texttt{move}, \texttt{release}, \texttt{cancel},
\texttt{step} et \texttt{snapshot}. Ouvrir le
\href{../../notebooks/module-4.ipynb}{notebook élève}, avec la copie de
correction conservée par l'enseignant.

\begin{Shaded}
\begin{Highlighting}[]
\ImportTok{export} \KeywordTok{const}\NormalTok{ inertiaOptions }\OperatorTok{=}\NormalTok{ \{ }\DataTypeTok{damping}\OperatorTok{:} \FloatTok{0.65}\OperatorTok{,} \DataTypeTok{restitution}\OperatorTok{:} \FloatTok{0.9}\NormalTok{ \}}\OperatorTok{;}
\end{Highlighting}
\end{Shaded}

La position est normalisée entre 0 et 1 ; la vitesse utilise ces unités
par seconde ; \texttt{dt} est une durée en secondes. À vitesse constante
de 0,2 unité/seconde pendant 0,05 seconde, le déplacement serait 0,01
unité. Ce petit calcul sert à comprendre les unités avant
l'amortissement.

Le code utilise \texttt{decay\ =\ Math.exp(-damping\ *\ dt)}. Si
\texttt{k} désigne l'amortissement et \texttt{t} une durée, la vitesse
devient \texttt{v(t)\ =\ v(0)\ ×\ exp(-k\ ×\ t)} entre les rebonds. Pour
\texttt{k\ \textgreater{}\ 0}, le déplacement correspondant est
\texttt{v(0)\ ×\ (1\ -\ exp(-k\ ×\ t))\ /\allowbreak{}\ k}. Le code traite aussi
\texttt{k\ =\ 0}. Ce modèle déterministe est un effet d'interface, sans
validation physique annoncée.

\section{Deux heures d'apprentissage
guidé}\label{deux-heures-dapprentissage-guiduxe9-3}

{\def\LTcaptype{none} % do not increment counter
\begin{longtable}[]{@{}
  >{\raggedright\arraybackslash}p{(\linewidth - 4\tabcolsep) * \real{0.3333}}
  >{\raggedright\arraybackslash}p{(\linewidth - 4\tabcolsep) * \real{0.3333}}
  >{\raggedright\arraybackslash}p{(\linewidth - 4\tabcolsep) * \real{0.3333}}@{}}
\toprule\noalign{}
\begin{minipage}[b]{\linewidth}\raggedright
Temps
\end{minipage} & \begin{minipage}[b]{\linewidth}\raggedright
Travail avec l'assistant
\end{minipage} & \begin{minipage}[b]{\linewidth}\raggedright
Trace sélectionnée
\end{minipage} \\
\midrule\noalign{}
\endhead
\bottomrule\noalign{}
\endlastfoot
0--30 min & Associer \texttt{x/\allowbreak{}y}, \texttt{vx/\allowbreak{}vy} et \texttt{dt} à leurs
unités. Faire un petit déplacement à vitesse constante. & Calcul court
et unités \\
30--60 min & Lire comment deux positions et deux instants donnent une
vitesse pendant le geste. Comparer \texttt{release} et \texttt{cancel}.
& Une séquence de données \\
60--90 min & Lire l'amortissement, le plafond de \texttt{dt} et les
rebonds. Expliquer pourquoi le survol ne réagrippe pas. & Une limite et
un scénario \\
90--120 min & Préparer la comparaison de 0,65 et 1,3 avec restitution
conservée à 0,9. & Prédiction et conditions \\
\end{longtable}
}

Demande possible : « Pars d'une position et d'une vitesse simples.
Explique pourquoi on multiplie par une durée en secondes. Fais-moi
vérifier l'unité de chaque valeur avant de comparer les amortissements.
» Demander ensuite à l'assistant de relier son exemple à \texttt{step},
sans créer un second moteur parallèle.

\section{Activité Colab --- 30
minutes}\label{activituxe9-colab-30-minutes-3}

\begin{enumerate}
\def\labelenumi{\arabic{enumi}.}
\tightlist
\item
  \textbf{Prédire --- 5 min.} Annoncer quel amortissement devrait
  raccourcir la trajectoire après relâchement. Conserver la zone, la
  restitution et un geste aussi comparable que possible.
\item
  \textbf{Modifier et observer --- 15 min.} Dans
  \texttt{student\_\allowbreak{}files}, remplacer uniquement \texttt{damping:\ 0.65}
  par \texttt{damping:\ 1.3}. Exécuter l'aperçu et comparer deux
  lancers. Croiser l'objet avec le pointeur, sans cliquer. Tester aussi
  un arrêt prolongé avant relâchement : dans ce modèle, plus de 0,12
  seconde sans nouvel échantillon annule l'élan estimé. Consigner ces
  gestes séparément.
\item
  \textbf{Expliquer et exporter --- 10 min.} Décrire la trajectoire, la
  variable changée et la comparabilité imparfaite des gestes. Remplir
  aide reçue, explication et limites, puis exporter le code modifié.
\end{enumerate}

Piste après tentative : un amortissement supérieur réduit plus
rapidement la vitesse. Deux lancers humains ne garantissent pas une
vitesse initiale identique. Pour isoler exactement l'effet du moteur,
les tests déterministes utilisent des entrées contrôlées ; votre
observation d'interface indique ce que vous avez effectivement essayé.

\section{Bilan solo --- 30 minutes}\label{bilan-solo-30-minutes-3}

Sélectionner les essais pendant 10 minutes, expliquer l'effet pendant 10
minutes et préparer une question pendant 10 minutes. Votre bilan doit
distinguer observation visuelle, calcul prévu et test déterministe. Une
cadence de rendu variable peut affecter l'expérience ; elle ne
transforme pas un écart en preuve de la théorie physique.

Conserver les paramètres et les unités est indispensable pour rejouer un
résultat. Une vidéo sans paramètres reste une trace partielle.

\section{Transfert et critères
observables}\label{transfert-et-crituxe8res-observables-3}

Adapter l'inertie à une carte HTML. L'élève définit les unités, conserve
des conditions comparables, montre pourquoi le pointeur passif ne change
pas la vitesse et décrit ce qui se fige en mode statique. Le modèle
garde la position courante et sa vitesse pendant le gel ; le temps
suspendu n'est pas à rattraper.

Les \href{../../../../../tests/learning-notebooks-bricks.test.ts}{tests
de briques} comparent notamment des durées contrôlées et des états
figés. Leur résultat courant appartient au bilan de validation.

\textbf{Ressources et statut :}
\href{https://developer.mozilla.org/en-US/docs/Web/API/Window/requestAnimationFrame}{requestAnimationFrame
--- MDN} ·
\href{../../../../../packages/learning/src/catalog.ts}{Mission du
catalogue} · \href{../../PROJECT.md}{Règles du projet et de la remise}

\section{Commandes de l'aperçu du Module
4}\label{commandes-de-laperuxe7u-du-module-4}

Le notebook corrigé conserve deux gestes distincts : les flèches
déplacent par pas de 0,04 et gardent \texttt{vx/\allowbreak{}vy} à zéro pendant le
geste; leur relâchement termine le geste sans lancer. Le
glisser-relâcher au pointeur estime une vitesse et produit l'inertie.
Espace fige/reprend l'image et le temps; Échap annule un geste actif. En
vol libre, le gel conserve position et vitesse, et sa durée n'est pas
rattrapée. Un geste actif est terminé sans lancement lors du gel.

Ces consignes suivent l'aperçu préparé dans le
\href{../../notebooks/module-4.ipynb}{notebook du Module 4}. Le reçu
antérieur qui avait détecté l'absence du clavier reste historique; la
qualification du correctif se lit dans la section commune des preuves.

\chapter{Module 5 --- Interpolation et
élasticité}\label{module-5-interpolation-et-uxe9lasticituxe9}

\textbf{Objectif : distinguer une transition vers une cible d'un
mouvement oscillant.} Vous allez relier cible, position, vitesse,
ressort et amortissement, puis expliquer un dépassement. Le module
occupe les jours 9--10 : \textbf{3 h solo, puis 1 h avec l'enseignant}.
Le projet personnel commence aussi pendant ces jours avec des heures
distinctes : consulter \href{../../PROJECT.md}{PROJECT.md}.

Le prérequis est d'avoir rencontré vitesse et durée. Le modèle de ce
module rend une transition d'interface réglable ; il ne revendique pas
une simulation physique validée.

\section{Intuition et code}\label{intuition-et-code}

L'\href{../../notebooks/module-5.ipynb}{activité Colab} utilise
\href{../../frontend/module-5/transitions.js}{transitions.js}. La copie
de correction conservée par l'enseignant reste séparée.

Une interpolation rapproche une valeur de sa cible selon une règle. Un
ressort conserve aussi une vitesse : même en arrivant à la cible, il
peut continuer et la dépasser. Une animation peut paraître douce avec
plusieurs mécanismes ; son apparence ne suffit pas à identifier le code.

\begin{Shaded}
\begin{Highlighting}[]
\ImportTok{export} \KeywordTok{const}\NormalTok{ transitionOptions }\OperatorTok{=}\NormalTok{ \{ }\DataTypeTok{stiffness}\OperatorTok{:} \DecValTok{36}\OperatorTok{,} \DataTypeTok{damping}\OperatorTok{:} \DecValTok{12}\NormalTok{ \}}\OperatorTok{;}
\end{Highlighting}
\end{Shaded}

Ici, \texttt{value} est une échelle sans unité, \texttt{target} sa cible
et \texttt{velocity} la variation de cette échelle par seconde. Le terme
\texttt{stiffness\ ×\ (target\ -\ value)} attire vers la cible ;
\texttt{damping\ ×\ velocity} freine. \texttt{step} borne la durée et la
découpe pour l'intégration. Les bornes de ce modèle sont explicites ;
des paramètres arbitraires ne deviennent pas sûrs parce qu'une seule
animation semble fonctionner.

\section{Deux heures d'apprentissage
guidé}\label{deux-heures-dapprentissage-guiduxe9-4}

{\def\LTcaptype{none} % do not increment counter
\begin{longtable}[]{@{}
  >{\raggedright\arraybackslash}p{(\linewidth - 4\tabcolsep) * \real{0.3333}}
  >{\raggedright\arraybackslash}p{(\linewidth - 4\tabcolsep) * \real{0.3333}}
  >{\raggedright\arraybackslash}p{(\linewidth - 4\tabcolsep) * \real{0.3333}}@{}}
\toprule\noalign{}
\begin{minipage}[b]{\linewidth}\raggedright
Temps
\end{minipage} & \begin{minipage}[b]{\linewidth}\raggedright
Travail avec l'assistant
\end{minipage} & \begin{minipage}[b]{\linewidth}\raggedright
Trace sélectionnée
\end{minipage} \\
\midrule\noalign{}
\endhead
\bottomrule\noalign{}
\endlastfoot
0--30 min & Dessiner une approche directe et un mouvement qui dépasse.
Nommer valeur et cible. & Deux schémas et une différence \\
30--60 min & Lire \texttt{set\allowbreak{}Target}, \texttt{velocity} et
\texttt{step}. Suivre un calcul d'accélération avec des nombres simples.
& Une ligne de calcul commentée \\
60--90 min & Lire amortissement, sous-pas et mode statique. Distinguer
figer la valeur et la remettre au départ. & Un scénario de
gel/reprise \\
90--120 min & Préparer la comparaison de 12 et 5 avec raideur 36 et
cible conservées. & Prédiction, variable et conditions \\
\end{longtable}
}

Demande possible : « Explique un ressort à partir de
\texttt{target\ -\ value} et \texttt{velocity}, puis montre ce qui
manque dans une simple interpolation. Demande-moi de prédire un
dépassement avant de l'essayer. » Une solution complète est permise ;
votre explication revient ensuite au code et à votre observation.

\section{Activité Colab --- 30
minutes}\label{activituxe9-colab-30-minutes-4}

\begin{enumerate}
\def\labelenumi{\arabic{enumi}.}
\tightlist
\item
  \textbf{Prédire --- 5 min.} Annoncer l'effet attendu du passage de
  \texttt{damping:\ 12} à \texttt{damping:\ 5}, avec la même cible et
  \texttt{stiffness:\ 36}. Prédire aussi ce que signifie activer le mode
  figé pendant le mouvement.
\item
  \textbf{Modifier et observer --- 15 min.} Modifier uniquement
  l'amortissement dans \texttt{student\_\allowbreak{}files}, puis exécuter l'aperçu.
  Comparer dépassement et stabilisation. Activer le gel à un instant
  choisi et examiner la valeur courante. La reprendre sans demander au
  modèle de rattraper le temps suspendu. Restaurer la valeur initiale
  pour un dernier essai.
\item
  \textbf{Expliquer et exporter --- 10 min.} Écrire ce que vous avez
  effectivement vu, votre aide reçue et une limite. Exporter la version
  choisie ; préciser si vous avez conservé ou restauré le paramètre.
\end{enumerate}

Piste après tentative : une diminution de l'amortissement permet
davantage d'oscillation dans cet essai. \texttt{step(dt,\ true)}
conserve la valeur et la vitesse présentes. Cela ne signifie pas rendre
le mouvement invisible ; l'image actuelle reste consultable.

\section{Bilan solo --- 30 minutes}\label{bilan-solo-30-minutes-4}

Passer 10 minutes à sélectionner les résultats, 10 minutes à expliquer
un dépassement et 10 minutes à préparer une question pour la revue.
Distinguer « j'ai observé cette oscillation » de « j'ai compris toutes
les conditions de stabilité ». Une incertitude conservée aide
l'enseignant à choisir l'exemple suivant.

Le projet personnel avance en parallèle ; ses heures ne sont pas prises
sur ce bilan de module.

\section{Transfert et critères
observables}\label{transfert-et-crituxe8res-observables-4}

Appliquer la transition à l'ouverture d'une fiche. L'élève nomme le
paramètre changé, distingue dépassement et stabilisation, conserve un
mode figé utilisable et montre où \texttt{dispose()} arrête la
ressource. Il explique la différence entre valeur cible et valeur
présente dans un nouvel exemple.

Les \href{../../../../../tests/learning-notebooks-bricks.test.ts}{tests
de briques} examinent les bornes et certains états statiques. Un test
réussi ne dispense pas d'observer le geste et de reformuler le
mécanisme.

\textbf{Ressources et statut :}
\href{https://developer.mozilla.org/en-US/docs/Web/CSS/@media/prefers-reduced-motion}{Préférence
de mouvement réduit --- MDN} ·
\href{../../../../../packages/learning/src/catalog.ts}{Mission et
critères du catalogue} · \href{../../PROJECT.md}{Projet personnel}

\chapter{Module 6 --- Particules et
ressources}\label{module-6-particules-et-ressources}

\textbf{Objectif : distinguer population, durée de vie, calcul et
rendu.} Vous allez modifier une population bornée, comparer les effets
et expliquer les ressources à libérer. Le module occupe les jours 11--12
: \textbf{3 h solo, puis 1 h avec l'enseignant}, avec \textbf{2 h solo
de projet supplémentaires et distinctes}.

Le prérequis est de lire une position, une vitesse et une durée. Une
particule de cet exercice est un petit ensemble de données ; elle n'est
pas une particule physique mesurée.

\section{Ce que contient une
particule}\label{ce-que-contient-une-particule}

Ouvrir
\href{../../frontend/module-6/particle-layer.js}{particle-layer.js}, le
\href{../../notebooks/module-6.ipynb}{notebook élève} et, pour
l'enseignant seulement, la copie de correction conservée par
l'enseignant.

\begin{Shaded}
\begin{Highlighting}[]
\ImportTok{export} \KeywordTok{const}\NormalTok{ particleOptions }\OperatorTok{=}\NormalTok{ \{ }\DataTypeTok{count}\OperatorTok{:} \DecValTok{64}\OperatorTok{,} \DataTypeTok{lifetime}\OperatorTok{:} \DecValTok{3}\OperatorTok{,} \DataTypeTok{seed}\OperatorTok{:} \DecValTok{17}\NormalTok{ \}}\OperatorTok{;}
\end{Highlighting}
\end{Shaded}

Chaque élément contient position \texttt{x/\allowbreak{}y}, vitesse \texttt{vx/\allowbreak{}vy} et
âge. La population est créée une fois. Un élément arrivé en fin de vie
réutilise son emplacement avec un âge nul et une nouvelle position. La
limite du modèle est de 256 éléments ; la durée de vie est bornée entre
0,2 et 10 secondes.

La graine \texttt{seed} rend la suite de nombres simulés répétable pour
les mêmes conditions. Elle ne prouve ni la validité physique du modèle
ni une performance identique sur tous les appareils. La fonction
\texttt{dispose()} vide les données de la couche ; le contrôleur de
scène libère séparément ses ressources graphiques.

\section{Deux heures d'apprentissage
guidé}\label{deux-heures-dapprentissage-guiduxe9-5}

{\def\LTcaptype{none} % do not increment counter
\begin{longtable}[]{@{}
  >{\raggedright\arraybackslash}p{(\linewidth - 4\tabcolsep) * \real{0.3333}}
  >{\raggedright\arraybackslash}p{(\linewidth - 4\tabcolsep) * \real{0.3333}}
  >{\raggedright\arraybackslash}p{(\linewidth - 4\tabcolsep) * \real{0.3333}}@{}}
\toprule\noalign{}
\begin{minipage}[b]{\linewidth}\raggedright
Temps
\end{minipage} & \begin{minipage}[b]{\linewidth}\raggedright
Travail avec l'assistant
\end{minipage} & \begin{minipage}[b]{\linewidth}\raggedright
Trace sélectionnée
\end{minipage} \\
\midrule\noalign{}
\endhead
\bottomrule\noalign{}
\endlastfoot
0--30 min & Lire les données d'une particule. Distinguer nombre, âge et
durée de vie. & Une particule annotée \\
30--60 min & Lire création et réutilisation. Expliquer pourquoi la
population n'augmente pas à chaque image. & Un parcours de fin de vie \\
60--90 min & Distinguer calcul JavaScript, dessin du navigateur, GPU et
runtime Python. Lire le nettoyage Three.js. & Une mesure pertinente et
une mesure hors portée \\
90--120 min & Préparer l'essai 64/128 avec durée 3 et graine 17
conservées. Choisir le scénario et l'appareil à décrire. & Prédiction et
conditions \\
\end{longtable}
}

Demande possible : « Suis une particule avant et après sa fin de vie.
Explique ce qui change et ce qui reste constant. Aide-moi ensuite à
choisir une mesure du navigateur, sans utiliser le temps Python comme
mesure du rendu Three.js. »

\section{Activité Colab --- 30
minutes}\label{activituxe9-colab-30-minutes-5}

\begin{enumerate}
\def\labelenumi{\arabic{enumi}.}
\tightlist
\item
  \textbf{Prédire --- 5 min.} Annoncer ce qui augmente lorsque
  \texttt{count} passe de 64 à 128. Prédire ce que vous pourriez voir et
  ce qu'un simple chronomètre du runtime ne permettrait pas de conclure.
\item
  \textbf{Modifier et observer --- 15 min.} Dans
  \texttt{student\_\allowbreak{}files}, changer seulement \texttt{count:\ 64} en
  \texttt{count:\ 128}. Garder \texttt{lifetime:\ 3} et
  \texttt{seed:\ 17}. Comparer la population et le comportement visuel.
  Activer le mode figé puis reprendre. Restaurer 64 après l'essai. Le
  notebook dessine un aperçu Canvas : il n'est pas une mesure de coût
  GPU Three.js.
\item
  \textbf{Expliquer et exporter --- 10 min.} Noter navigateur, scénario,
  population et durée. Si une métrique manque, écrire « indisponible »
  plutôt que zéro. Conserver observation, aide reçue et limite, puis
  exporter la version restaurée.
\end{enumerate}

Piste après tentative : deux fois plus d'éléments signifie plus de
données à mettre à jour et à dessiner. Cela ne garantit pas une durée de
frame exactement doublée : appareil, navigateur, autres travaux et
stratégie de rendu comptent. Une mesure de performance doit décrire sa
source et ses conditions.

\section{Bilan solo --- 30 minutes}\label{bilan-solo-30-minutes-5}

Choisir les traces pendant 10 minutes, expliquer population et
ressources pendant 10 minutes, puis préparer une question pendant 10
minutes. Le rendu peut inclure une mesure absente ou un appareil peu
performant : ces limites sont utiles pour choisir un effet facultatif.

Vous devez pouvoir distinguer « 128 éléments dans les données » de « 128
objets graphiques indépendants ». Le renderer peut organiser les données
autrement ; il faut lire son chemin de code avant de généraliser.

\section{Transfert et critères
observables}\label{transfert-et-crituxe8res-observables-5}

Créer une couche facultative, bornée et désactivable qui signale une
sélection. L'élève distingue population et durée, indique la source
d'une mesure, conserve un libellé sans effet et montre le nettoyage à la
fermeture. Une population identique après plusieurs cycles ne suffit pas
à prouver toute l'absence de fuite ; le parcours complet doit être
observé.

Le code garde les données présentes en mode figé. La couche peut être
désactivée explicitement pour une préférence d'accessibilité distincte.

\textbf{Ressources et statut :}
\href{https://threejs.org/manual/\#en/cleanup}{Nettoyage Three.js ---
manuel officiel} ·
\href{../../frontend/module-3/scene-controller.js}{Contrôleur de scène
et ressources graphiques} ·
\href{../../../../../packages/learning/src/catalog.ts}{Mission du
catalogue} · \href{../../PROJECT.md}{Projet et conditions de remise}

\chapter{Module 7 --- Parcours, états et
mémoire}\label{module-7-parcours-uxe9tats-et-muxe9moire}

\textbf{Objectif : gérer les transitions sans état bloqué.} Vous allez
conserver une sélection, limiter un historique et expliquer la
restauration d'une vue. Le module occupe les jours 13--14 : \textbf{3 h
solo, puis 1 h avec l'enseignant}, avec \textbf{2 h solo de projet
distinctes}.

Le Module 1 gérait un geste ; celui-ci gère le parcours. « Quel pointeur
est actif ? » et « Quelle fiche est ouverte ? » appartiennent à des
états différents, même lorsque le même clic intervient dans les deux.

\section{Code du parcours}\label{code-du-parcours}

Ouvrir
\href{../../frontend/module-7/interaction-flow.js}{interaction-flow.js},
le \href{../../notebooks/module-7.ipynb}{notebook élève} et, pour
l'enseignant, la copie de correction conservée par l'enseignant.

\begin{Shaded}
\begin{Highlighting}[]
\ImportTok{export} \KeywordTok{const}\NormalTok{ flowOptions }\OperatorTok{=}\NormalTok{ \{ }\DataTypeTok{historyLimit}\OperatorTok{:} \DecValTok{32}\NormalTok{ \}}\OperatorTok{;}
\end{Highlighting}
\end{Shaded}

Le modèle conserve \texttt{view}, \texttt{selection}, \texttt{elapsed}
et \texttt{static}. Les vues autorisées sont \texttt{mission},
\texttt{experience}, \texttt{proofs} et \texttt{summary}. La sélection
conserve un identifiant borné. \texttt{navigate} ajoute une entrée ; si
le budget est dépassé, la plus ancienne sort.

L'aperçu Colab reçoit \texttt{host:\ null}. Son historique est interne :
il ne pilote pas le bouton Retour de la page Colab. Dans le projet
Astro, l'hôte navigateur relie les paramètres \texttt{learning\allowbreak{}View} et
\texttt{learning\allowbreak{}Selection} à l'URL et écoute \texttt{popstate}. Les deux
parcours sont vérifiés séparément.

\section{Deux heures d'apprentissage
guidé}\label{deux-heures-dapprentissage-guiduxe9-6}

{\def\LTcaptype{none} % do not increment counter
\begin{longtable}[]{@{}
  >{\raggedright\arraybackslash}p{(\linewidth - 4\tabcolsep) * \real{0.3333}}
  >{\raggedright\arraybackslash}p{(\linewidth - 4\tabcolsep) * \real{0.3333}}
  >{\raggedright\arraybackslash}p{(\linewidth - 4\tabcolsep) * \real{0.3333}}@{}}
\toprule\noalign{}
\begin{minipage}[b]{\linewidth}\raggedright
Temps
\end{minipage} & \begin{minipage}[b]{\linewidth}\raggedright
Travail avec l'assistant
\end{minipage} & \begin{minipage}[b]{\linewidth}\raggedright
Trace sélectionnée
\end{minipage} \\
\midrule\noalign{}
\endhead
\bottomrule\noalign{}
\endlastfoot
0--30 min & Dessiner liste → fiche → retour avec une sélection précise.
Nommer ce qui doit être conservé. & Trois états du parcours \\
30--60 min & Lire \texttt{navigate} et \texttt{snapshot}. Repérer
validation, ajout et retrait d'une entrée. & Une trace de cinq
changements \\
60--90 min & Lire URL, \texttt{popstate} et visibilité. Distinguer
historique interne et navigation navigateur. & Une condition de
restauration \\
90--120 min & Préparer l'essai des limites 32/4 et le gel de la
chronologie. & Prédiction et scénario \\
\end{longtable}
}

Demande possible : « Aide-moi à reproduire un retour à une mauvaise
sélection. Montre quelles données restaurer, puis demande-moi de tracer
cinq changements avec un historique limité à quatre. » L'assistant
propose le scénario ; vous conservez la séquence effectivement utilisée.

\section{Activité Colab --- 30
minutes}\label{activituxe9-colab-30-minutes-6}

\begin{enumerate}
\def\labelenumi{\arabic{enumi}.}
\tightlist
\item
  \textbf{Prédire --- 5 min.} Écrire quelle entrée devrait sortir après
  cinq changements lorsque \texttt{history\allowbreak{}Limit} vaut 4. Choisir une
  séquence précise, avec au moins deux éléments sélectionnés. Prédire
  l'effet du mode figé sur \texttt{elapsed}.
\item
  \textbf{Modifier et observer --- 15 min.} Dans
  \texttt{student\_\allowbreak{}files}, comparer \texttt{history\allowbreak{}Limit:\ 32} puis
  \texttt{history\allowbreak{}Limit:\ 4} avec la même séquence. Consulter les entrées
  conservées et la sélection courante. Activer le gel, attendre,
  reprendre et examiner l'horloge. Ne pas appeler cet essai une
  validation du bouton Retour navigateur.
\item
  \textbf{Expliquer et exporter --- 10 min.} Décrire l'entrée retirée,
  l'état restauré et les conditions de chronologie. Noter aide, limite
  et question, puis exporter le code effectivement choisi.
\end{enumerate}

Piste après tentative : l'entrée la plus ancienne sort lorsque la
cinquième est ajoutée à un budget de quatre. Le budget concerne ici les
entrées de diagnostic internes, pas une suppression arbitraire de
l'historique global du navigateur. Le gel suspend \texttt{elapsed} et ne
doit pas ajouter le temps d'attente au retour.

\section{Bilan solo --- 30 minutes}\label{bilan-solo-30-minutes-6}

Prendre 10 minutes pour sélectionner la séquence, 10 minutes pour
expliquer un retour et 10 minutes pour préparer le webinaire. Le rendu
conserve état initial, événements et état final. « Cela semble
fonctionner » devient une vérification utile lorsqu'on peut retrouver le
scénario précis.

Une trace locale n'est pas automatiquement un journal à publier. L'élève
choisit les éléments à remettre ; aucune conversation complète n'est
nécessaire.

\section{Transfert et critères
observables}\label{transfert-et-crituxe8res-observables-6}

Créer un parcours fiche → liste qui restaure la sélection. L'élève
conserve vue et identifiant, explique l'entrée retirée par la borne,
distingue historique interne et historique navigateur et vérifie le gel
sans rattrapage de temps. Après fermeture, les événements doivent être
nettoyés, pas ajoutés à nouveau sans limite.

Les \href{../../../../../tests/learning-notebooks-bricks.test.ts}{tests
logiciels de briques} examinent notamment l'état interne. Le parcours
navigateur demande une observation dans l'environnement cible.

\textbf{Ressources et statut :}
\href{https://developer.mozilla.org/en-US/docs/Web/API/History_API}{History
API --- MDN} ·
\href{https://developer.mozilla.org/en-US/docs/Web/API/Page_Visibility_API}{Page
Visibility API --- MDN} ·
\href{../../../../../packages/learning/src/catalog.ts}{Mission du
catalogue} · \href{../../PROJECT.md}{Projet et reprise}

\section{Commandes de l'aperçu du Module
7}\label{commandes-de-laperuxe7u-du-module-7}

Les boutons des trois éléments affichent \texttt{element-1},
\texttt{element-2} et \texttt{element-3}. Tab puis Entrée ou Espace
permettent de les choisir; le pointeur déclenche le même choix. Le
snapshot affiche \texttt{selection} et les boutons signalent la
sélection courante. Changer ensuite \texttt{mission},
\texttt{experience}, \texttt{proofs} ou \texttt{summary} conserve cet
identifiant.

Préparer cinq choix précis, avec au moins deux identifiants différents,
puis répéter la même séquence après avoir changé \texttt{history\allowbreak{}Limit}
de 32 à 4. Les choix d'élément et les changements de vue ajoutent chacun
une entrée. Noter l'entrée retirée et la sélection finale. Cet
historique reste celui de l'iframe; le bouton Retour du navigateur est
examiné dans le projet Astro. Le statut du retest du
\href{../../notebooks/module-7.ipynb}{notebook corrigé} se lit dans la
qualification commune.

\chapter{Module 8 --- Assistant, preuves et
capacités}\label{module-8-assistant-preuves-et-capacituxe9s}

\textbf{Objectif : comprendre ce qu'un agent peut lire, calculer et
proposer.} Vous allez examiner permissions, révisions et provenance,
puis découvrir comment raccorder vos productions. Le module occupe les
jours 15--16 : \textbf{3 h solo, puis 1 h avec l'enseignant}. La
dernière heure solo du projet et son heure de clôture sont
\textbf{supplémentaires et distinctes}.

Les états, le rendu et le parcours ont été rencontrés. Une capacité
d'agent ajoute une question : quel contenu est effectivement partagé,
pour quelle révision et avec quelle autorité ?

\section{Code, trace et appel réel}\label{code-trace-et-appel-ruxe9el}

Ouvrir
\href{../../frontend/module-8/register-capabilities.js}{register-capabilities.js}
et le \href{../../notebooks/module-8.ipynb}{notebook élève}. La copie de
correction est réservée à l'enseignant.

Le module prépare \textbf{un outil du projet élève},
\texttt{student\_\allowbreak{}get\_\allowbreak{}learning\_\allowbreak{}snapshot}. Il ne remplace pas le
registre des vingt-cinq outils d'Orbit. Son adaptateur reçoit
\texttt{can\allowbreak{}Read}, \texttt{get\allowbreak{}Revision}, \texttt{get\allowbreak{}Snapshot} et l'API
navigateur. Il contrôle la permission avant de lire et vérifie la
révision. Les résultats distinguent notamment
\texttt{CONSENT\_\allowbreak{}REQUIRED}, \texttt{STALE\_\allowbreak{}REVISION},
\texttt{UNAVAILABLE} et \texttt{READY}.

\begin{Shaded}
\begin{Highlighting}[]
\ControlFlowTok{if}\NormalTok{ (}\OperatorTok{!}\FunctionTok{canRead}\NormalTok{()) }\ControlFlowTok{return}\NormalTok{ \{ }\DataTypeTok{state}\OperatorTok{:} \StringTok{\textquotesingle{}CONSENT\_REQUIRED\textquotesingle{}}\NormalTok{ \}}\OperatorTok{;}
\end{Highlighting}
\end{Shaded}

La présence de \texttt{READY} dans une trace préparée n'est pas la
preuve d'un appel réel. Un assistant ne peut pas fabriquer une
approbation humaine en ajoutant un champ. Le texte d'une source, d'un
export ou d'une proposition est une donnée à examiner, pas une
instruction qui change les permissions.

\section{Deux heures d'apprentissage
guidé}\label{deux-heures-dapprentissage-guiduxe9-7}

{\def\LTcaptype{none} % do not increment counter
\begin{longtable}[]{@{}
  >{\raggedright\arraybackslash}p{(\linewidth - 4\tabcolsep) * \real{0.3333}}
  >{\raggedright\arraybackslash}p{(\linewidth - 4\tabcolsep) * \real{0.3333}}
  >{\raggedright\arraybackslash}p{(\linewidth - 4\tabcolsep) * \real{0.3333}}@{}}
\toprule\noalign{}
\begin{minipage}[b]{\linewidth}\raggedright
Temps
\end{minipage} & \begin{minipage}[b]{\linewidth}\raggedright
Travail avec l'assistant
\end{minipage} & \begin{minipage}[b]{\linewidth}\raggedright
Trace sélectionnée
\end{minipage} \\
\midrule\noalign{}
\endhead
\bottomrule\noalign{}
\endlastfoot
0--30 min & Lire la capacité bornée. Identifier contenu autorisé,
permission et révision. & Une fiche de capacité \\
30--60 min & Examiner la trace préparée du notebook. Repérer une
permission absente, une version périmée ou une autorité usurpée. &
Objection reliée à un champ \\
60--90 min & Lire les outils d'Orbit, passage, portée et proposition.
Distinguer présence d'une citation et pertinence. & Un exemple de limite
de conclusion \\
90--120 min & Préparer un besoin utile pour votre assistant et la
comparaison d'une limite de volume. & Prédiction et périmètre choisi \\
\end{longtable}
}

Demande possible : « Explique les conditions de lecture dans ce code.
Ensuite, examine cette trace comme une donnée d'exercice : quelles
affirmations ne puis-je pas vérifier à partir d'elle ? Propose une
capacité utile et bornée pour mon projet. » Votre assistant peut
expliquer directement ; vous conservez la justification de la permission
et du contenu partagé.

\section{Activité Colab --- 30
minutes}\label{activituxe9-colab-30-minutes-7}

\begin{enumerate}
\def\labelenumi{\arabic{enumi}.}
\tightlist
\item
  \textbf{Prédire --- 5 min.} Écrire pourquoi une trace préparée
  \texttt{READY} ne certifie ni un appel natif ni une revue humaine.
  Prédire l'effet d'une limite de snapshot passant de 12 000 à 6 000
  caractères de sérialisation.
\item
  \textbf{Modifier et observer --- 15 min.} Dans
  \texttt{student\_\allowbreak{}files}, modifier seulement la comparaison de taille
  \texttt{\textgreater{}\ 12000} en \texttt{\textgreater{}\ 6000}.
  Conserver \texttt{can\allowbreak{}Read}, \texttt{expected\allowbreak{}Revision}, les
  revérifications et le nettoyage. Examiner la trace préparée ; relier
  vos objections aux champs de permission et de révision. Ce carnet
  \textbf{ne réalise pas un appel WebMCP natif} et ne réimplémente aucun
  moteur en Python.
\item
  \textbf{Expliquer et exporter --- 10 min.} Décrire ce que votre
  capacité lirait, la limite retenue et l'autorité de sa réponse.
  Conserver aide reçue, observation, limite et question, puis exporter
  votre version.
\end{enumerate}

Piste après tentative : la borne concerne ici la longueur d'une chaîne
JSON, pas un nombre de tokens de modèle ou une limite universelle
d'octets. Réduire le volume ne doit pas supprimer les contrôles de
permission ou de version. Après changement de révision, la capacité doit
être enregistrée pour la nouvelle révision.

\section{Bilan solo --- 30 minutes}\label{bilan-solo-30-minutes-7}

Consacrer 10 minutes au choix des preuves, 10 minutes à votre
explication et 10 minutes à préparer la revue. Distinguer trace
d'exercice, proposition de l'assistant, calcul d'un moteur et décision
humaine. Une affirmation peut être incertaine sans que cela attribue une
note à l'élève.

\section{Relier les productions à la fin de la
leçon}\label{relier-les-productions-uxe0-la-fin-de-la-leuxe7on}

Vous avez conservé huit briques. \textbf{Elles peuvent maintenant former
un socle de frontend :} geste et état, carte, scène, inertie,
transition, particules, parcours et capacité.

Pendant les 25 minutes de construction du webinaire du Module 8, dans la
partie finale du notebook 8, charger vos \textbf{huit ZIP réels}. Le
raccordement fourni vérifie les manifestes et prépare un projet avec vos
fichiers. Il ne comble pas une absence avec un corrigé caché. Si un
module manque ou si une empreinte ne correspond pas, conserver l'erreur
et réexporter le bon rendu.

Le code d'intégration est fourni et attribué au cours dans
\href{../../assembly/README.md}{assembly/}. Vos modifications sont dans
les dossiers \texttt{src/\allowbreak{}learning/\allowbreak{}module-N/\allowbreak{}}. L'assemblage du ZIP ne
compile pas Astro : la compilation et le parcours du projet sont
vérifiés séparément. Une modification hors du contrat appelle une
adaptation expliquée ; elle ne justifie pas une substitution silencieuse
du travail.

\section{Transfert et critères
observables}\label{transfert-et-crituxe8res-observables-7}

Préparer une capacité de lecture utile à votre projet, limitée à une
sélection explicitement partagée. L'élève distingue trace et appel réel,
retrouve provenance et portée, conserve la révision et explique la
différence proposition/décision. Il montre une modification personnelle
dans le projet assemblé.

Dans Orbit, les moteurs \texttt{baseline}, \texttt{n} et \texttt{p}
restent des évaluations séparées sur les mêmes données choisies. T/I/F
classe les preuves d'une affirmation ; il \textbf{ne note pas la
compréhension de l'élève}. Aucun vote ni score moyen ne choisit un
moteur automatiquement.

\textbf{Ressources et statut :}
\href{https://developer.chrome.com/docs/ai/webmcp/imperative-api}{API
impérative WebMCP --- Chrome} ·
\href{https://www.sanity.io/docs/ai/sanity-context}{Sanity Context ---
documentation officielle} ·
\href{../../../../../packages/learning/src/catalog.ts}{Contrats et
mission du catalogue} · \href{../../assembly/README.md}{Assemblage
fourni} et \href{../../PROJECT.md}{projet personnel}

\chapter{Projet personnel --- adapter, intégrer et
expliquer}\label{projet-personnel-adapter-intuxe9grer-et-expliquer}

\textbf{1 h avec l'enseignant pour cadrer → 6 h solo → 1 h avec
l'enseignant pour clôturer.} L'élève choisit son intention et adapte les
productions qu'il a réellement conservées.

\section{Calendrier --- seize jours, sans double
comptage}\label{calendrier-seize-jours-sans-double-comptage}

{\def\LTcaptype{none} % do not increment counter
\begin{longtable}[]{@{}
  >{\raggedright\arraybackslash}p{(\linewidth - 4\tabcolsep) * \real{0.3333}}
  >{\raggedright\arraybackslash}p{(\linewidth - 4\tabcolsep) * \real{0.3333}}
  >{\raggedright\arraybackslash}p{(\linewidth - 4\tabcolsep) * \real{0.3333}}@{}}
\toprule\noalign{}
\begin{minipage}[b]{\linewidth}\raggedright
Jours
\end{minipage} & \begin{minipage}[b]{\linewidth}\raggedright
Module : 3 h solo + 1 h avec l'enseignant
\end{minipage} & \begin{minipage}[b]{\linewidth}\raggedright
Projet : heures supplémentaires
\end{minipage} \\
\midrule\noalign{}
\endhead
\bottomrule\noalign{}
\endlastfoot
1--2 & Module 1 & --- \\
3--4 & Module 2 & --- \\
5--6 & Module 3 & --- \\
7--8 & Module 4 & --- \\
9--10 & Module 5 & 1 h de cadrage avec l'enseignant + 1 h solo \\
11--12 & Module 6 & 2 h solo \\
13--14 & Module 7 & 2 h solo \\
15--16 & Module 8 & 1 h solo + 1 h de clôture avec l'enseignant \\
\textbf{Projet} & & \textbf{6 h solo + 2 h avec l'enseignant} \\
\end{longtable}
}

Le projet commence avec les notions déjà rencontrées. Son intégration
finale bénéficie de la découverte du Module 8 ; le thème personnel et
ses premiers choix peuvent être explorés avant cette découverte.

\section{Cadrage du projet --- 1 heure avec
l'enseignant}\label{cadrage-du-projet-1-heure-avec-lenseignant}

L'idée appartient à l'élève. Le cadrage transforme cette idée en besoin
observable et en périmètre réalisable. L'assistant propose des pistes ;
l'élève conserve la décision et sa justification.

{\def\LTcaptype{none} % do not increment counter
\begin{longtable}[]{@{}
  >{\raggedright\arraybackslash}p{(\linewidth - 4\tabcolsep) * \real{0.3333}}
  >{\raggedright\arraybackslash}p{(\linewidth - 4\tabcolsep) * \real{0.3333}}
  >{\raggedright\arraybackslash}p{(\linewidth - 4\tabcolsep) * \real{0.3333}}@{}}
\toprule\noalign{}
\begin{minipage}[b]{\linewidth}\raggedright
Durée
\end{minipage} & \begin{minipage}[b]{\linewidth}\raggedright
Action
\end{minipage} & \begin{minipage}[b]{\linewidth}\raggedright
Production
\end{minipage} \\
\midrule\noalign{}
\endhead
\bottomrule\noalign{}
\endlastfoot
10 min & Décrire la personne, son besoin et l'action principale &
Intention en trois phrases \\
15 min & Comparer deux pistes et choisir un parcours court & Choix
motivé et limite \\
20 min & Définir contenu, interaction, contribution 3D, alternative HTML
et capacité utile & Carte du projet et critères \\
15 min & Choisir une expérience, une remise et les points de revue &
Plan solo et trace initiale \\
\end{longtable}
}

Exemples possibles : une galerie de réalisations, un explorateur de
trois notions, une présentation interactive d'un projet ou un guide
visuel. Ce sont des pistes ; elles ne remplacent pas l'intention
personnelle.

La contribution Three.js doit servir un usage expliqué. Le contenu
conserve un parcours HTML et clavier. L'élève peut modifier la
présentation ou désactiver un effet facultatif, avec une justification.
Les scènes du cours sont indépendantes du landing d'Orbit.

\section{Travail autonome du projet --- 6
heures}\label{travail-autonome-du-projet-6-heures}

{\def\LTcaptype{none} % do not increment counter
\begin{longtable}[]{@{}
  >{\raggedright\arraybackslash}p{(\linewidth - 4\tabcolsep) * \real{0.3333}}
  >{\raggedright\arraybackslash}p{(\linewidth - 4\tabcolsep) * \real{0.3333}}
  >{\raggedright\arraybackslash}p{(\linewidth - 4\tabcolsep) * \real{0.3333}}@{}}
\toprule\noalign{}
\begin{minipage}[b]{\linewidth}\raggedright
Bloc
\end{minipage} & \begin{minipage}[b]{\linewidth}\raggedright
Travail avec l'assistant
\end{minipage} & \begin{minipage}[b]{\linewidth}\raggedright
Trace à préparer
\end{minipage} \\
\midrule\noalign{}
\endhead
\bottomrule\noalign{}
\endlastfoot
\textbf{1 h --- jours 9--10} & Clarifier intention et contenu, puis
préparer une première carte et une interaction & Besoin, choix,
modification et question \\
\textbf{2 h --- jours 11--12} & Adapter la présentation et la scène ;
vérifier lisibilité, sélection et alternative HTML & Fichiers modifiés,
essai, limites \\
\textbf{2 h --- jours 13--14} & Raccorder le parcours ; mener une
expérience comparative ; préparer permissions et reprise & Paramètres,
prédiction, observations et restauration \\
\textbf{1 h --- jours 15--16} & Intégrer les huit exports, examiner le
résultat et préparer l'explication finale & Projet, instructions de
reprise et bilan sélectionné \\
\end{longtable}
}

Pour chaque bloc, choisir une aide appropriée : brainstorm, exploration,
enseignement, guidage, débogage, critique, transfert ou révision. Le
journal sélectionné décrit l'aide reçue et la décision de l'élève. Il
n'exige pas l'historique complet d'une conversation personnelle.

La découverte du Module 8 utilise
\href{../../assembly/README.md}{l'assemblage fourni} et
\href{../../assembly/merge_modules.py}{merge\_modules.py}. Le ZIP
construit importe les huit dossiers dans \texttt{src/\allowbreak{}learning/\allowbreak{}}.
\href{../../assembly/src/student-frontend.js}{student-frontend.js} est
le raccordement fourni, avec une seule boucle de rendu et une fermeture
des ressources. L'élève doit pouvoir retrouver au moins une modification
de chacun de ses rendus dans le projet assemblé.

Les dépendances directes du
\href{../../assembly/package.json}{package.json} sont fixées. Le
lockfile produit lors de la compilation vérifiée doit accompagner une
reproduction des dépendances transitives. \textbf{Créer un ZIP dans
Colab ne compile pas Astro.} Une compilation réussie doit être
enregistrée dans son environnement réel ; un aperçu seul ne la prouve
pas.

\section{Remise sélectionnée}\label{remise-suxe9lectionnuxe9e}

Le rendu final contient : le projet ; les huit ZIP de module ; une
intention ; une expérience comparative ; une explication personnelle ;
l'aide déclarée ; des limites ; les instructions de reprise. Les données
nécessaires à un essai restent bornées au périmètre choisi.

Chaque export de module comprend les briques réellement présentes dans
\texttt{student\_\allowbreak{}files}, un résultat JSON, une copie rejouable du
notebook, \texttt{INTEGRATION.md} et un manifeste d'empreintes. Le
résultat indique \texttt{external-declared}. L'empreinte identifie un
contenu ; elle ne prouve ni une exécution indépendante, ni l'auteur de
chaque ligne, ni la compréhension.

Un import ne transforme pas automatiquement ce résultat en approbation
humaine. L'enseignant peut refaire l'essai, demander une reformulation
ou proposer un transfert. Un rendu périmé est conservé avec sa révision
; il ne remplace pas le travail courant.

\section{Critères observables --- une grille de
revue}\label{crituxe8res-observables-une-grille-de-revue}

{\def\LTcaptype{none} % do not increment counter
\begin{longtable}[]{@{}
  >{\raggedright\arraybackslash}p{(\linewidth - 4\tabcolsep) * \real{0.3333}}
  >{\raggedright\arraybackslash}p{(\linewidth - 4\tabcolsep) * \real{0.3333}}
  >{\raggedright\arraybackslash}p{(\linewidth - 4\tabcolsep) * \real{0.3333}}@{}}
\toprule\noalign{}
\begin{minipage}[b]{\linewidth}\raggedright
Critère
\end{minipage} & \begin{minipage}[b]{\linewidth}\raggedright
Ce que l'élève peut montrer
\end{minipage} & \begin{minipage}[b]{\linewidth}\raggedright
Ce qui reste distinct
\end{minipage} \\
\midrule\noalign{}
\endhead
\bottomrule\noalign{}
\endlastfoot
Intention & Une personne, un besoin et un parcours explicites & Une
préférence esthétique \\
Attribution & Code fourni, modifié, aide reçue et raccordement
identifiés & Une signature ou une empreinte \\
État & Geste, sélection, annulation et retour expliqués dans le code &
Une capture finale seule \\
Accessibilité & Libellés, clavier, contenu HTML et mode figé utilisables
& La présence d'un bouton de préférence \\
Expérience & Prédiction, variable, conditions, observation et
restauration & Un score ou une impression sans protocole \\
Capacité d'agent & Lecture bornée, partage et révision montrés & Une
approbation humaine automatique \\
Reprise & Fichiers, dépendances, export et instructions examinables &
Une compilation supposée \\
Transfert & Une notion appliquée dans un contexte nouveau & Une
définition mémorisée \\
\end{longtable}
}

Un critère peut être \textbf{rencontré, pratiqué, démontré sur un
scénario ou encore non examiné}. Ces catégories ne déclarent pas une
maîtrise universelle. L'enseignant garde sa décision et sa justification
séparées des suggestions de l'assistant.

\section{Stockage et permissions}\label{stockage-et-permissions}

{\def\LTcaptype{none} % do not increment counter
\begin{longtable}[]{@{}
  >{\raggedright\arraybackslash}p{(\linewidth - 4\tabcolsep) * \real{0.3333}}
  >{\raggedright\arraybackslash}p{(\linewidth - 4\tabcolsep) * \real{0.3333}}
  >{\raggedright\arraybackslash}p{(\linewidth - 4\tabcolsep) * \real{0.3333}}@{}}
\toprule\noalign{}
\begin{minipage}[b]{\linewidth}\raggedright
Destination
\end{minipage} & \begin{minipage}[b]{\linewidth}\raggedright
Ce qu'elle peut contenir
\end{minipage} & \begin{minipage}[b]{\linewidth}\raggedright
Choix de l'élève
\end{minipage} \\
\midrule\noalign{}
\endhead
\bottomrule\noalign{}
\endlastfoot
Orbit local & Dossier, journal et artefacts sélectionnés & Autoriser la
sauvegarde ou rester en mémoire et exporter \\
Colab / Drive personnel & Copie du carnet, code, texte et sorties
conservées & Choisir ce qui est partagé et avec qui \\
Sanity & Contenu explicitement accepté comme publiable & Examiner
contenu et destination avant écriture \\
\end{longtable}
}

Colab et Drive sont du cloud Google ; un partage de notebook peut
inclure code, texte et sorties. Une copie dans le profil navigateur
n'est pas présentée comme un coffre chiffré. Le cours n'exige aucun
secret, appel d'API payante, GPU ou abonnement Colab Pro. Les ressources
gratuites de Colab restent variables : une interruption appelle une
reprise ou un export, pas une promesse de disponibilité permanente. Voir
la \href{https://research.google.com/colaboratory/faq.html}{FAQ
officielle Colab}.

Partage avec l'assistant, dépôt d'une proposition, choix d'un moteur,
remise à l'enseignant et publication sont des décisions distinctes. Le
plugin du Studio de l'élève utilise ses propres droits. Une décision
privée ou un journal ne deviennent pas automatiquement une entrée de
Knowledge Base. Les instructions contenues dans une source restent des
données.

Les trois moteurs d'Orbit restent séparés sur des entrées choisies :
\texttt{baseline}, \texttt{n} et \texttt{p}. Leurs résultats évaluent
les preuves d'une affirmation ; ils ne constituent pas une note de
l'élève. Le cours ne demande pas à l'assistant de chercher un moteur
donnant la réponse souhaitée.

\chapter{Qualification technique et lecture des
preuves}\label{qualification-technique-et-lecture-des-preuves}

Édition pédagogique : \texttt{orbit-course-1.0.0}. Point documentaire du
\textbf{1er octobre 2026}. Les résultats ci-dessous décrivent les
périmètres des reçus consultés et des observations du coordinateur ; la
composition du manuel ne relance aucun test. Le
\href{../../../../receipts/FORMATION_DELIVERY_STATUS.md}{reçu de
livraison maintenu} conserve sa chronologie ; les reçus précis liés ici
gouvernent cette édition lorsque ce résumé antérieur n'est pas encore
synchronisé.

\textbf{État de livraison : \texttt{delivery\allowbreak{}Ready:\ false},
\texttt{full\allowbreak{}Mission\allowbreak{}Complete:\ false}.} Le site de formation et les
logiciels ont leurs contrôles réussis ; le nouvel assemblage des huit
exports a également passé ses 35 contrôles dans un vrai navigateur. Le
transport public du cours Sanity Context reste \texttt{HOLD} et les
pages Canva ne sont pas encore produites. La session du second Studio a
été examinée en lecture seule ; ses appels d'agent et écritures restent
distincts. Cette édition documente les mécanismes utilisables et leurs
limites ; elle ne déclare pas l'ensemble de la mission terminé.

{\def\LTcaptype{none} % do not increment counter
\begin{longtable}[]{@{}
  >{\raggedright\arraybackslash}p{(\linewidth - 4\tabcolsep) * \real{0.3333}}
  >{\raggedright\arraybackslash}p{(\linewidth - 4\tabcolsep) * \real{0.3333}}
  >{\raggedright\arraybackslash}p{(\linewidth - 4\tabcolsep) * \real{0.3333}}@{}}
\toprule\noalign{}
\begin{minipage}[b]{\linewidth}\raggedright
Objet examiné
\end{minipage} & \begin{minipage}[b]{\linewidth}\raggedright
Résultat rapporté dans le reçu consulté
\end{minipage} & \begin{minipage}[b]{\linewidth}\raggedright
Portée de ce résultat
\end{minipage} \\
\midrule\noalign{}
\endhead
\bottomrule\noalign{}
\endlastfoot
Contrats et régressions logiciels historiques & Kaggle :
\href{../../../../receipts/formation/software-133-final.json}{133/133
assertions sur le payload corrigé} ; le reçu conserve le premier run de
117 et les tests Python séparés & Payload \texttt{cae9a0f4…} uniquement
; les 117 historiques appartiennent à la version
\href{https://www.kaggle.com/code/celebrum/orbit-formation-software-validation?scriptVersionId=354372770}{354372770} \\
Checkpoints logiciels historiques conservés &
\href{../../../../receipts/formation/software-139-kaggle.json}{139/139
assertions}, puis
\href{../../../../receipts/formation/software-141-kaggle.json}{141/141
assertions, exit 0 et Python exit 0}, source \texttt{28813ecb…f72531} &
Exécutions interactives historiques du draft Kaggle ; chaque payload
conserve sa source. La sauvegarde immuable et les étapes ultérieures
gardent leur reçu propre \\
Checkpoint logiciel historique avant v5 &
\href{../../../../receipts/formation/software-141-final-kaggle.json}{141/141
assertions, exit 0 et logiciel notebook Python exit 0}, source
\texttt{98a0c659…62ed87} & Ce reçu décrit un draft antérieur. Il ne
reçoit pas rétroactivement les empreintes ou les résultats de la version
publique courante \\
Version publique Kaggle historique v4 &
\href{../../../../receipts/formation/software-version4-kaggle.json}{Readback
conservé} de
\href{https://www.kaggle.com/code/celebrum/orbit-formation-software-validation?scriptVersionId=354422952}{la
version 4, identifiant 354422952}, durée 105,3 s : 141 tests dans 15
suites, sept tests Python, compilation des huit exports indépendants,
second Studio statique verrouillé et 29 tests PHP/185 assertions & Zéro
appel de modèle. Le payload source \texttt{98a0c659…} et les sorties
gardent cette exécution historique ; ils ne deviennent pas ceux de v5 \\
Version publique Kaggle courante v5 &
\href{../../../../receipts/formation/software-version5-kaggle.json}{Readback
conservé} de
\href{https://www.kaggle.com/code/celebrum/orbit-formation-software-validation?scriptVersionId=354436854}{la
version 5, identifiant 354436854}, durée 120,7 s : 141 tests dans 15
suites, sept tests Python, nouvel assemblage exact \texttt{141466603…},
second Studio statique verrouillé et 29 tests PHP/185 assertions &
Payload source \texttt{68949498…e7966e}, zéro appel de modèle. Ce succès
logiciel ne valide ni l'admission publique Context, ni une compréhension
humaine, ni les écritures du Studio. Les sorties compilées gardent leur
propre exécution et empreinte \\
Formation publique et outils natifs & Navigateur E2B piloté par Kaggle :
\href{../../../../receipts/formation/live-browser-72.json}{72/72
contrôles réussis} & Release exercée, fixtures synthétiques, zéro appel
de modèle ; découverte de 25 outils, appels choisis, stockage, archive
et versions \\
Formation publique historique après correction du manifeste &
\href{../../../../receipts/formation/live-browser-final.json}{72/72
contrôles réussis}, run
\texttt{formation-790dfa8e-614e-4646-85d0-e428a65307fa}, durée 5 147 ms,
sur le domaine public & 25 outils natifs, permissions, huit expériences,
sauvegarde, versions et ZIP exercés. Le manifeste public contient 358
fichiers ; HTML, JS, CSS et archive portable f2 correspondent aux
empreintes de cette exécution. La marque, les cinq portes du landing, le
guide et le petit écran sont contrôlés. Fixtures techniques, zéro appel
de modèle \\
Formation publique courante après le raccordement M3/M8 &
\href{../../../../receipts/formation/live-browser-pointer-target-final.json}{Nouveau
parcours 72/72 réussi}, run
\texttt{formation-de718cf4-ef62-4333-a2ac-aadf1f49a267}, durée 5 093 ms
& Navigateur E2B piloté par Kaggle sur le domaine public, zéro appel de
modèle et aucune erreur runtime. Les scénarios restent des fixtures
techniques ; la compréhension et la revue humaine ne sont pas déduites
de ce succès \\
Sources Colab --- QA groupée historique & Huit exports et huit
relectures dans des espaces Python séparés & Un même runtime CPU ; une
relecture de namespace diffère d'un redémarrage réel \\
Aperçus Colab --- sondes historiques & Huit rapports de sondes
navigateur & Vérifications par scripts ; trace du Module 8 préparée et
gestes pointeur de confiance encore à qualifier \\
Assemblage des huit vrais exports & Liaison de neuf fichiers de briques,
46 fichiers de manifeste vérifiés, compilation Astro réussie dans Kaggle
& Build statique avec Node 22.20.0, Astro 7.3.3 et Three.js 0.181.2 ;
comportement navigateur du projet assemblé distinct \\
Reprise après redémarrage Colab réel &
\href{../../../../receipts/formation/colab-recovery.json}{Un redémarrage
suivi de la relecture des sources capturées} & Chemin de checkpoint
public synthétique ; anciens rapports navigateur importés plutôt que
rejoués ; copies froides et reprise privée distinctes \\
Corrections historiques de stockage, intégrité ZIP et versions & Les
résultats historiques 133/133 et 72/72 incluent ces parcours dans leurs
scénarios & Le lecteur ZIP ensuite partagé avec le Studio possède sa
qualification ultérieure dans v5 et le reçu portable f2 ; le premier
checkpoint reste distinct \\
Installation statique dans un second Studio &
\href{../../../../receipts/formation/second-studio-static.json}{PASS\_\allowbreak{}STATIC\_\allowbreak{}INSTALL\_\allowbreak{}BUILD}
pour l'archive \texttt{3e3a499d…} & Installation, import et build
statique réussis ; UI authentifiée, changement de workspace et transport
Context encore distincts \\
Installation avec dépendances figées après lecteur ZIP partagé &
\href{../../../../receipts/formation/portable-studio-locked-install-kaggle.json}{PASS\_\allowbreak{}STATIC\_\allowbreak{}INSTALL\_\allowbreak{}BUILD},
archive \texttt{ea0f93ef…992b90} & Installation par \texttt{npm\ ci},
import et build statique observés dans le draft Kaggle ; session
authentifiée, Context cross-origin, WebMCP natif et écritures Sanity
restent \texttt{false} \\
Installation de l'archive portable courante &
\href{../../../../receipts/formation/portable-studio-v2-install-kaggle.json}{PASS\_\allowbreak{}STATIC\_\allowbreak{}INSTALL\_\allowbreak{}BUILD},
archive \texttt{f2b0dc68…f9e973} & \texttt{npm\ ci}, import et build
statique dans un hôte vierge observés ; 315 fichiers. Authentification,
Context cross-origin, WebMCP natif, écriture Sanity et publication npm
restent \texttt{false} \\
Session du second Studio authentifié &
\href{../../../../receipts/formation/authenticated-second-studio.json}{Lecture
DOM réelle} sur \texttt{/\allowbreak{}formation/\allowbreak{}studio/\allowbreak{}orbit-learning-lab} :
utilisateur authentifié visible, huit modules, FR/EN/ES, stockage
mémoire et six permissions désactivées & L'UI indique le registre
enregistré. Aucun appel d'agent natif ni écriture Content Lake n'est
exécuté dans ce reçu ; le changement de workspace et la publication
sélectionnée ne reçoivent pas de validation implicite \\
Copies Colab indépendantes &
\href{../../../../receipts/formation/free-colab-module-1-independent.json}{Module
1},
\href{../../../../receipts/formation/free-colab-module-2-independent.json}{Module
2},
\href{../../../../receipts/formation/free-colab-module-3-independent.json}{Module
3} : copies neuves, modification et export observés & Fixtures opérées
par Codex ; incidents conservés, CPU observé et forfait gratuit rapporté
par l'utilisateur ; compréhension humaine encore à examiner \\
Incident puis retest clavier du Module 4 &
\href{../../../../receipts/formation/free-colab-module-4-independent.json}{Le
reçu antérieur} détecte le défaut clavier ;
\href{../../../../receipts/formation/free-colab-module-4-independent-retest.json}{le
retest corrigé} observe flèches, relâchement à vitesse nulle, lancer au
pointeur, rebond et gel/reprise dans une copie CPU neuve du commit
\texttt{784dbca} & Export effectivement téléchargé, SHA-256
\texttt{7e6c739c…4d67a5}, runtime arrêté. Annulation tactile native et
comparaison quantitative d'amortissement restent non vérifiées ;
compréhension non examinée \\
Copies indépendantes des Modules 5, 6 et 8 &
\href{../../../../receipts/formation/free-colab-module-5-independent.json}{Module
5},
\href{../../../../receipts/formation/free-colab-module-6-independent.json}{Module
6} et
\href{../../../../receipts/formation/free-colab-module-8-independent.json}{Module
8} : modification, export reçu et arrêt du runtime observés & M5 observe
une oscillation après amortissement modifié ; M6 confirme population et
commande statique, sans mesure des pixels/âges ; M8 conserve une trace
préparée, sans appel WebMCP natif ni exécution de sa limite de
snapshot \\
Incident puis retest des sélections du Module 7 &
\href{../../../../receipts/formation/free-colab-module-7-independent.json}{Le
reçu indépendant antérieur} garde
\texttt{selection\allowbreak{}Variants\allowbreak{}Verified:\ false} ;
\href{../../../../receipts/formation/free-colab-module-7-independent-retest.json}{le
retest corrigé} observe choix clavier d'\texttt{element-2} et
d'\texttt{element-3}, conservation entre vues, budget 32/4 et gel
d'\texttt{elapsed} & Copie CPU neuve du commit \texttt{4b06da41…4d1bc6},
export reçu \texttt{82752841…97f3ce}, runtime arrêté. Le véritable
bouton Retour navigateur et la compréhension restent non examinés \\
Import manuel des huit exports dans le Module 8 & Le
\href{../../../../receipts/formation/free-colab-module-8-independent.json}{reçu
M8} observe le champ multiple et deux délais d'attente du dialogue de
fichiers, sans fichier transmis ; l'annulation produit le refus attendu
faute de huit exports & \texttt{manual\allowbreak{}Upload\allowbreak{}Verified:\ false}. Le
transport HTTP de huit archives synthétiques choisies et leur nouvel
assemblage gardent leurs propres preuves ; ils ne valident pas le
dialogue manuel \\
Assemblage des huit exports indépendants &
\href{../../../../receipts/formation/free-colab-module-8-independent-assembly.json}{Assemblage
réellement exécuté dans Colab CPU}, huit archives exactes téléchargées
par HTTP, empreintes/tailles/identifiants contrôlés, neuf fichiers de
briques et 46 fichiers de manifeste & Archive \texttt{96acc0dd…5b2aed},
78 147 octets, export reçu et runtime arrêté.
\texttt{ASSEMBLED\_\allowbreak{}NOT\_\allowbreak{}BUILT} côté Colab ; fixtures de QA, pas
production d'un élève \\
Compilation des exports indépendants &
\href{../../../../receipts/formation/actual-independent-assembly-static-kaggle.json}{Compilation
des vrais exports réussie} dans Kaggle : Node 22.20.0, Astro 7.3.3,
Three.js 0.181.2, \texttt{npm\ ci} avec lock fourni, build exit 0 & Les
huit modules et neuf briques exportées sont importés et appelés, sans
remplacement par un corrigé ; sortie \texttt{71982b43…}. Comportement
navigateur, WebMCP natif et compréhension restent \texttt{false} dans ce
reçu \\
Premier parcours natif de l'assemblage indépendant &
\href{../../../../receipts/formation/actual-assembly-browser-first-failure.json}{Incident
initial conservé}, puis
\href{../../../../receipts/formation/actual-assembly-browser-keyboard-failure.json}{dix
contrôles réussis avant l'échec clavier} sur le vrai artefact compilé
\texttt{96acc0dd…} & Les empreintes du HTML/JS servi, la scène, le
consentement et la capacité native bornée sont observés.
\texttt{keyboard-card-selects-real-observation} échoue ;
\texttt{success:\ false}, \texttt{completed:\ false}. Le diagnostic et
le retest ne reçoivent aucun succès anticipé \\
Défaut réel de sélection 3D dans l'assemblage précédent & Le
\href{../../../../receipts/formation/actual-assembly-browser-raycast-failure.json}{parcours
avec le geste clavier corrigé} atteint quinze contrôles réussis, puis
échoue sur \texttt{actual-raycast-selects-second-sphere} & La capture du
Module 1 dirige le relâchement vers le stage, alors que le Module 3
l'écoutait sur son canvas. La compilation de \texttt{96acc0dd…} ne
démontrait pas cette intégration. L'incident reste conservé après le
correctif de la cible d'événements \\
Nouvelle copie indépendante du Module 3 &
\href{../../../../receipts/formation/free-colab-module-3-pointer-target-independent.json}{Reprise
CPU du correctif}, source \texttt{6a60234}, caméra 5 puis 7, sélection
HTML au clavier et clics 3D observés & Export \texttt{b202088c…7c3335},
runtime arrêté. La cible par défaut du canvas est exercée ;
\texttt{shared\allowbreak{}Stage\allowbreak{}Integration\allowbreak{}Verified:\ false} dans ce reçu. La taille
projetée n'a pas été mesurée et la compréhension n'a pas été examinée \\
Nouvelle copie indépendante du Module 8 &
\href{../../../../receipts/formation/free-colab-module-8-pointer-target-independent.json}{Nouvelle
copie CPU} du raccordement corrigé et export effectivement reçu & Export
\texttt{8d80562e…bd1654}. La trace du notebook reste préparée ; cet
export ne démontre pas un appel natif WebMCP \\
Nouvel assemblage des exports corrigés &
\href{../../../../receipts/formation/free-colab-module-8-pointer-target-assembly.json}{Assemblage
CPU réellement reçu} : huit archives par HTTP avec contrôles conservés,
six exports inchangés, M3/M8 fraîchement exportés, neuf briques et 46
fichiers de manifeste & Archive \texttt{14146660…b0b1ba}, 79 149 octets,
\texttt{ASSEMBLED\_\allowbreak{}NOT\_\allowbreak{}BUILT} côté Colab. Le dialogue manuel, la
compilation Kaggle et la sélection sous capture du stage conservent
chacun leur validation propre \\
Compilation du nouvel assemblage exact &
\href{../../../../receipts/formation/actual-pointer-target-assembly-static-kaggle.json}{Build
Kaggle réussi} pour \texttt{141466603…} : huit modules, neuf briques, 46
payloads vérifiés, \texttt{npm\ ci} et Astro build exit 0 & Node
22.20.0, Astro 7.3.3, Three.js 0.181.2, lock \texttt{423e4878…a0edd8}.
Fichiers originaux conservés, aucune brique remplacée par un corrigé. Le
succès statique complète le reçu Colab sans modifier son statut
historique \\
Navigateur du nouvel assemblage exact &
\href{../../../../receipts/formation/actual-pointer-target-assembly-browser-kaggle.json}{35/35
contrôles réels réussis}, run
\texttt{assembly-038e527c-85d9-44d8-bc93-0997ccb771e7}, durée 4 917 ms &
Kaggle pilote E2B sur l'artefact compilé \texttt{141466603…}. Sélection
HTML au clavier, raycasting sous capture du stage, inertie, gel,
particules, repli sans WebGL, capacité native bornée et nettoyage
exercés ; zéro erreur runtime et zéro appel de modèle. Fixtures de QA,
pas travail d'élève ni approbation humaine \\
Régressions du serveur public et privé &
\href{../../../../receipts/formation/php-context-29-pass-kaggle.json}{PASS\_\allowbreak{}PHP\_\allowbreak{}REGRESSIONS},
exit 0 ; le coordinateur observe 29 tests et 185 assertions sur le
payload \texttt{5fb59cbe…} & PHP 8.5.11 du runtime FrankenPHP officiel
v1.12.7 ; routes, format ciblé, cours, KB, comptes et espaces de
travail. Fixtures sans credentials réels, activation publique de
production \texttt{false} dans ce run \\
Navigateur de la release candidate &
\href{../../../../receipts/formation/candidate-browser-62-v2.json}{62/62
contrôles formation} et
\href{../../../../receipts/formation/atom-five-navigation-candidate.json}{77/77
contrôles atome} & Kaggle pilote E2B contre le serveur candidat
\texttt{127.0.0.1:4321}, avec zéro appel de modèle ; ces résultats ne
constituent pas la vérification de la release publique \\
Atome et cinq destinations sur le domaine public &
\href{../../../../receipts/formation/atom-five-navigation-live.json}{77/77
contrôles de l'atome} sur \texttt{https:/\allowbreak{}/\allowbreak{}orbit.securedme.ca/\allowbreak{}} &
Navigateur E2B piloté par Kaggle ; zéro appel de modèle. Ce parcours
artistique et ses liens gardent leur périmètre distinct des mécanismes
de formation et de la validité scientifique \\
Incidents publics du manifeste conservés &
\href{../../../../receipts/formation/live-browser-v2-digest-failure.json}{Premier
échec du digest} puis
\href{../../../../receipts/formation/live-browser-latest-digest-failure.json}{échec
de la lecture courante} : la copie compressée du manifeste restait
ancienne alors que l'archive distante était f2 & Le coordinateur a
identifié une réponse Brotli ancienne distincte de la réponse sans
compression, puis corrigé les dates et la compression du manifeste
uniquement. Le nouveau reçu 72/72 établit le readback réussi ; les
incidents gardent leurs résultats initiaux \\
Critères encore ouverts & Transport public Context, appels natifs et
publication du second Studio, documents Canva & \textbf{Qualification de
mission encore incomplète} ; les versions Kaggle, le navigateur de
l'assemblage, le readback formation et la lecture du Studio authentifié
sont établis, mais ces résultats ne ferment pas les autres critères \\
Save \& Run All après lecteur ZIP partagé & Le coordinateur rapporte
pour
\href{https://www.kaggle.com/code/celebrum/orbit-formation-software-validation?scriptVersionId=354391847}{la
version Kaggle 354391847} : 133 assertions et assemblage des exports
réussis, puis échec \texttt{ETARGET} pour \texttt{motion-utils@\^{}13.5}
pendant l'installation du second Studio & \textbf{Exécution partielle} :
les étapes réussies ne valident pas le second Studio ni toute la
campagne ; son reçu immutable et son graphe exact restent à rattacher au
reçu maintenu \\
Graphe de dépendances de la reprise & Le lockfile préparatoire
\texttt{studio-success-lock.json}, SHA-256 \texttt{e4215188…412938},
conserve \texttt{motion-utils} résolu en \texttt{13.3.0}. Le nouveau
reçu statique conserve sa propre empreinte d'installation
\texttt{ed94c135…b4ba06} & Préparation et exécution ultérieure ont des
empreintes distinctes ; le succès de la reprise conserve l'échec initial
\texttt{ETARGET} \\
Transport du corpus pédagogique et accès serveur & GET
\texttt{/api/v1/course-context/\{outline,entries\}}, credentials omis,
neuf sources publiques prévues, empreintes et portée refusées en cas de
désaccord ;
\href{../../../../receipts/formation/php-context-29-pass-kaggle.json}{29
régressions PHP} réussies & Le serveur est construit pour refuser une
admission incomplète. Une réussite logicielle ne garantit pas que Sanity
a produit les neuf références attendues ; son activation publique reste
désactivée \\
Admission réelle de Sanity Context &
\href{../../../../receipts/formation/course-context-import-status.json}{Neuf
imports publics terminés, puis admission refusée}. La dernière
construction guidée est \texttt{succeeded}, mais
\texttt{learning\_\allowbreak{}modules} cite les Modules 1 à 7 et deux sources
anciennes ; le Module 8 est prêt sans être cité, \texttt{PROJECT.md}
reste \texttt{skipped}, et le namespace demandé n'est pas créé &
\texttt{public\allowbreak{}Admitted\allowbreak{}Entries:\ 0},
\texttt{public\allowbreak{}Endpoint\allowbreak{}Enabled:\ false},
\texttt{HOLD\_\allowbreak{}PUBLIC\_\allowbreak{}COURSE\_\allowbreak{}CONTEXT\_\allowbreak{}ADMISSION}. Aucune entrée,
provenance ou revue humaine n'est fabriquée. Aucun journal personnel ni
credential n'est inclus. Le serveur attend un corpus entièrement
admissible ; son déploiement et son cache de production ne sont pas
prouvés par le reçu PHP \\
Refus contrôlé du serveur réellement déployé &
\href{../../../../receipts/formation/course-context-production-disabled.json}{Transport
désactivé observé depuis Kaggle} : les deux routes du cours retournent
HTTP 503, \texttt{UNAVAILABLE}, \texttt{CONTEXT\_\allowbreak{}NOT\_\allowbreak{}READY} et
\texttt{no\allowbreak{}Fallback:\ true} & CORS public \texttt{*} sans credentials,
cache \texttt{no-store,\ private}, \texttt{nosniff}. Une origine
étrangère reçoit HTTP 403 sur l'ancienne route privée de KB. Archive
serveur \texttt{0246a362…}. Ce refus conforme ne devient pas une lecture
réussie du corpus Context \\
Propriétaire du registre natif & Source inspectée : registre par
document et propriétaire de montage, fermeture du prédécesseur, rejet du
nettoyage tardif d'une ancienne vue & Le code exprime la règle ; sa
nouvelle preuve native appartient aux contrôles de la version
courante \\
Retrait de l'accès temporaire de mission &
\href{../../../../receipts/formation/mission-token-retirement.json}{RETIRED\_AND\_REFUSED,
HTTP 401} & Credential dédié à la mission enregistré comme retiré ;
aucune valeur secrète dans ce manuel \\
Compréhension de l'élève & À examiner dans les rendus et les séances &
Reformulation, démonstration et transfert humains \\
\end{longtable}
}

Le reçu Colab groupé conserve \texttt{PARTIAL\_\allowbreak{}VALIDATION},
\texttt{cold\allowbreak{}Starts\allowbreak{}Verified:\ false},
\texttt{all\allowbreak{}Course\allowbreak{}Functions\allowbreak{}Validated:\ false} et
\texttt{learner\allowbreak{}Understanding:\ NOT\_\allowbreak{}EXAMINED}. Le redémarrage et la
reprise ultérieurs sont un reçu séparé : ils ne réécrivent pas les
champs du premier rapport. Le coordinateur rapporte une observation de
l'interface CPU d'un compte gratuit ; le JSON de cette exécution garde
\texttt{account\allowbreak{}Plan\allowbreak{}Observed:\ not-observed}. Ces provenances restent
distinctes.

Pour lire une preuve, nommer \textbf{la version, l'entrée,
l'environnement, le scénario, le résultat et sa limite}. Une empreinte
peut montrer qu'un fichier correspond à sa copie. Le statut
\texttt{external-declared} conserve ce qu'une personne déclare avoir
fait. Une exécution technique demande un résultat d'exécution ; un
examen pédagogique demande une explication ou un transfert réellement
examiné.

L'ensemble des campagnes Gemini, les trajectoires autonomes de modèles
et les résultats pédagogiques relèvent de leurs propres expériences. Les
contrôles de ce cours n'en déclarent pas l'achèvement.

Le résultat \texttt{PASS\_\allowbreak{}STATIC\_\allowbreak{}INSTALL\_\allowbreak{}BUILD} historique concerne
son archive et son graphe d'installation. Les reçus ultérieurs
établissent leurs nouvelles installations statiques par
\texttt{npm\ ci}, avec des empreintes distinctes. La version 354391847
conserve son échec \texttt{ETARGET} ; elle ne reçoit pas rétroactivement
le statut de cette reprise. L'authentification réelle, le transport
GET/CORS et une publication humaine sélectionnée gardent leurs propres
critères de sortie. Le build Kaggle de l'assemblage indépendant complète
le résultat Colab \texttt{ASSEMBLED\_\allowbreak{}NOT\_\allowbreak{}BUILT} sans modifier ses
champs, ni attribuer le succès au dialogue d'import manuel indisponible.

\section{Module 4 --- correctif et retest
observé}\label{module-4-correctif-et-retest-observuxe9}

La source élève corrigée a l'empreinte \texttt{d8ff5bd2…c2d2ce}; son
corrigé séparé a l'empreinte \texttt{3acb9425…66807b}. La brique
\texttt{inertia.js} conserve son contrat. Le changement porte sur
l'aperçu : flèches bornées avec vitesse nulle, fin de geste sans lancer,
inertie du pointeur, Espace et gestion des interruptions. Le reçu du
retest constate \texttt{x:\ 0.5\ →\ 0.54} et \texttt{y:\ 0.5\ →\ 0.54},
puis \texttt{held:\ false}, \texttt{vx/\allowbreak{}vy:\ 0} au relâchement. Le geste
pointeur lance réellement l'objet ; le rebond et le gel de
position/vitesse ont été observés. Son export téléchargé reste une
fixture de QA opérée par Codex, séparée du corrigé et d'une production
d'élève.

\section{Module 7 --- correctif et retest
observé}\label{module-7-correctif-et-retest-observuxe9}

Le notebook du Module 7 corrigé porte l'empreinte
\texttt{6b85b378…cc0f45}; son corrigé séparé porte
\texttt{1426f798…1b031b}. Le changement concerne l'aperçu HTML, avec les
identifiants \texttt{element-1}, \texttt{element-2} et
\texttt{element-3}. Le contrat de la brique et l'objectif de deux
sélections restent inchangés. Le retest utilise Entrée puis Espace pour
sélectionner les deuxième et troisième éléments, examine
\texttt{aria-pressed} et répète cinq choix avec les limites 32 puis 4.
La première entrée sort sous la borne de quatre; la vue et la sélection
finales restent \texttt{summary} et \texttt{element-3}. Le gel
d'\texttt{elapsed} est observé. Ces résultats concernent l'historique
interne de la copie Colab, distinct du véritable retour du projet Astro.

\section{Modules 3 et 8 --- intégration du geste et retest
observé}\label{modules-3-et-8-intuxe9gration-du-geste-et-retest-observuxe9}

Le défaut de l'ancien assemblage concernait l'intégration, malgré sa
compilation réussie : le Module 1 capturait le geste sur le stage,
tandis que le Module 3 attendait son relâchement sur le canvas.
\texttt{pointer\allowbreak{}Target} reste facultatif pour l'usage isolé et reçoit
explicitement le stage dans l'hôte d'assemblage. La projection et le
raycasting utilisent le rectangle du canvas. La copie CPU neuve du
Module 3 observe les caméras 5 puis 7, les choix HTML et les clics 3D ;
elle ne prétend pas, seule, valider le stage commun.

Les nouvelles archives des Modules 3 et 8 ont ensuite remplacé
uniquement leurs anciennes copies dans le manifeste de QA. Les six
autres exports sont conservés à l'identique. Le nouvel assemblage réel
\texttt{141466603…} est compilé dans Kaggle, puis exercé dans E2B avec
35 contrôles réussis. Ces reçus établissent la sélection sous capture
partagée sur ce scénario. Les échecs de \texttt{96acc0dd…} restent dans
la chronologie ; aucune vérification n'est attribuée rétroactivement à
cet artefact. L'empreinte du manifeste Git HTTP et celle de la copie
Windows CRLF sont distinctes et identifiées dans le reçu Colab ; les
assertions de contenu ont été conservées.

\section{Retour d'usage et prochaine
revue}\label{retour-dusage-et-prochaine-revue}

Le coordinateur transmet le retour de l'utilisateur : le laboratoire
fonctionne pour son essai et ses boutons sont compréhensibles. La revue
graphique reste une activité commune ultérieure. Ce retour borne un
usage humain ; les critères techniques se lisent dans leurs reçus. Les
documents en préparation poursuivent le périmètre pédagogique actuel ;
le prochain tutoriel et les éditoriaux restent dans leurs prochaines
missions.

Cette section est le point unique de mise à jour du manuel pour les
nouveaux reçus. Avant les exports définitifs, le coordinateur y reporte
les liens immuables, empreintes et limites du dernier payload, puis
recompose les sources. Une qualification technique réussie ouvre la
revue documentaire ; la disponibilité des PDF et des pages Canva reste
une autre étape. Le \href{../TOOLCHAIN_AND_VERSIONS.md}{registre des
versions} distingue versions déclarées, exécutions observées et outils
d'export encore à qualifier.

\section{Lire la suspension du corpus
public}\label{lire-la-suspension-du-corpus-public}

Les neuf documents de cours ont été importés. Deux tentatives guidées
utiles ont ensuite terminé leur construction ; leurs contenus ne
satisfont pas la portée stricte du cours. Une construction marquée
réussie chez le fournisseur décrit son exécution, pas l'admissibilité de
sa sortie pour Orbit. Le compteur conserve 107 documents utilisés, avec
une différence d'une unité signalée après la construction, sans ajout
public supposé au-delà des neuf imports. Aucun problème ancien de KB n'a
été accepté ou rejeté pour obtenir un succès.

Pour cette édition, l'assistant peut examiner les missions, les
artefacts partagés et les passages fournis dans le dossier. Il ne doit
pas présenter un sommaire Context public du cours comme disponible. Le
chemin désactivé retourne \texttt{UNAVAILABLE}, raison
\texttt{CONTEXT\_\allowbreak{}NOT\_\allowbreak{}READY}, avec \texttt{no\allowbreak{}Fallback:\ true} ; cette
réponse est validée par les régressions serveur et par le reçu de
production qui garde ses observations HTTP. La reprise exige des
références admissibles couvrant les huit modules et le protocole, puis
la validation d'une lecture du corpus effectivement admis.

\chapter{Partie II --- accompagner et
enseigner}\label{partie-ii-accompagner-et-enseigner}

Cette partie prépare l'enseignant à partir du travail sélectionné par
l'élève. Les corrigés donnent des pistes de raisonnement ; l'enseignant
compare ces pistes au scénario et à l'explication effectivement
présentés. Les scripts sont des propositions de parole pour
Jean-Sébastien, avec des pauses et intonations suggérées.

\section{Préparer une séance}\label{pruxe9parer-une-suxe9ance}

Avant le webinaire, choisir une observation, une difficulté et un
transfert. Ouvrir le fichier réellement remis et sa version. Repérer ce
qui a été fourni, modifié, déclaré exécuté et effectivement examiné.
Vérifier que la prédiction et les conditions permettent une comparaison
utile. Une limite déclarée peut devenir le meilleur exemple de la
séance.

Pendant la séance, montrer un effet, laisser une prédiction, relier une
donnée au code et faire essayer une modification. La question finale
change le contexte afin d'examiner le transfert. L'assistant reste
disponible pour expliquer ; l'élève garde la décision et montre ce qu'il
retient.

\section{Examiner sans confondre les
preuves}\label{examiner-sans-confondre-les-preuves}

{\def\LTcaptype{none} % do not increment counter
\begin{longtable}[]{@{}
  >{\raggedright\arraybackslash}p{(\linewidth - 4\tabcolsep) * \real{0.3333}}
  >{\raggedright\arraybackslash}p{(\linewidth - 4\tabcolsep) * \real{0.3333}}
  >{\raggedright\arraybackslash}p{(\linewidth - 4\tabcolsep) * \real{0.3333}}@{}}
\toprule\noalign{}
\begin{minipage}[b]{\linewidth}\raggedright
Trace
\end{minipage} & \begin{minipage}[b]{\linewidth}\raggedright
Question de revue
\end{minipage} & \begin{minipage}[b]{\linewidth}\raggedright
Portée
\end{minipage} \\
\midrule\noalign{}
\endhead
\bottomrule\noalign{}
\endlastfoot
Code & Où se trouve la modification ? Quelle aide a contribué ? &
Contenu et attribution déclarée \\
Observation & Quel geste, appareil et paramètre produisent ce résultat ?
& Scénario observé \\
Empreinte & Les contenus comparés sont-ils identiques ? & Identité des
fichiers \\
Test logiciel & Quelles entrées et quel comportement sont contrôlés ? &
Assertion technique bornée \\
Reformulation & L'élève explique-t-il le mécanisme avec ses données ? &
Explication examinée \\
Transfert & Utilise-t-il la notion dans une situation nouvelle ? & Usage
nouveau examiné \\
\end{longtable}
}

Choisir les statuts \textbf{rencontré, pratiqué, démontré sur ce
scénario, encore à examiner}. Une réussite locale ne devient pas une
maîtrise universelle. T/I/F porte sur les preuves d'une affirmation ; il
ne devient pas une note d'élève.

\section{Intervenir à partir d'une
difficulté}\label{intervenir-uxe0-partir-dune-difficultuxe9}

{\def\LTcaptype{none} % do not increment counter
\begin{longtable}[]{@{}
  >{\raggedright\arraybackslash}p{(\linewidth - 4\tabcolsep) * \real{0.3333}}
  >{\raggedright\arraybackslash}p{(\linewidth - 4\tabcolsep) * \real{0.3333}}
  >{\raggedright\arraybackslash}p{(\linewidth - 4\tabcolsep) * \real{0.3333}}@{}}
\toprule\noalign{}
\begin{minipage}[b]{\linewidth}\raggedright
Situation rencontrée
\end{minipage} & \begin{minipage}[b]{\linewidth}\raggedright
Intervention proposée
\end{minipage} & \begin{minipage}[b]{\linewidth}\raggedright
Trace examinable
\end{minipage} \\
\midrule\noalign{}
\endhead
\bottomrule\noalign{}
\endlastfoot
L'élève reconnaît l'effet mais hésite sur le code & Demander quel champ
change et ouvrir le symbole indiqué dans la fiche & Champ montré et
phrase personnelle \\
La prédiction et l'observation diffèrent & Relever paramètre, conditions
et scénario avant d'expliquer l'écart & Comparaison conservée avec sa
limite \\
L'aide a donné la solution complète & Demander une modification courte
ou un nouveau contexte & Décision de l'élève et essai de transfert \\
Une vérification est indisponible & Nommer la condition manquante et
choisir une reprise & Statut explicite et prochaine action \\
\end{longtable}
}

Pour chaque revue, conserver le scénario examiné, l'explication, la
limite, l'aide déclarée, le transfert et la suite choisie. Le retour de
l'assistant et la décision humaine restent séparés. Les fiches
détachables reprennent ce cadre pour pouvoir circuler sans le manuel
complet.

\section{Choisir l'aide}\label{choisir-laide}

L'élève peut demander un indice conceptuel, un indice technique, du
pseudo-code, un exemple partiel, une solution complète ou son
explication. Après une solution complète, proposer une vérification et
un changement de contexte. Le niveau d'aide reçu est déclaré avec la
production ; il ne sert pas à disqualifier automatiquement
l'apprentissage.

\section{Préparer l'accessibilité}\label{pruxe9parer-laccessibilituxe9}

Conserver la liste HTML, les libellés et le clavier. Le mode statique
garde l'image courante et suspend son évolution. L'alternative
accessible garde le contenu lorsque WebGL ou le réseau sont
indisponibles. Contraste, texte agrandi et langue améliorent le parcours
; le code personnel n'est pas traduit automatiquement.

\chapter{Fiche enseignant 1 --- Geste, état et
rendu}\label{fiche-enseignant-1-geste-uxe9tat-et-rendu}

\textbf{But de la revue :} L'élève relie l'événement à la position et
nomme la fin du pointeur actif ; l'enseignant observe son nouveau
scénario.

\textbf{Préparation :} ouvrir le rendu sélectionné, le notebook et
\href{../../frontend/module-1/interaction-state.js}{interaction-state.js}.
Le \href{../../notebooks/instructor/module-1.ipynb}{corrigé séparé} est
une piste à examiner après la tentative. Retrouver la prédiction et la
version réellement choisie avant de comparer.

\section{Une activité avec l'élève et son
assistant}\label{une-activituxe9-avec-luxe9luxe8ve-et-son-assistant}

Faire annoncer la position après deux flèches droites, puis exécuter les
deux pas clavier. L'élève retrouve ensuite pointerId dans le code et
propose un scénario d'interruption.

L'assistant peut expliquer ou fournir un exemple complet. L'élève
choisit ensuite sa modification, nomme l'aide reçue et montre un
résultat. La question à poser est : \textbf{« Quelle donnée change
pendant le geste et quelle condition autorise ce changement ? »}

Comparer \textbf{keyboardStep : 0.04 → 0.08}, en conservant
\textbf{zone, position initiale et nombre de pressions}. Le rendu
sélectionné peut prendre la forme de schéma événement → état → rendu,
code choisi et séquence observée. La limite à garder visible est :
\textbf{un Échap ou un événement scripté ne démontre pas une annulation
tactile native}.

\section{Déroulé et script de
séance}\label{duxe9rouluxe9-et-script-de-suxe9ance}

{\def\LTcaptype{none} % do not increment counter
\begin{longtable}[]{@{}
  >{\raggedright\arraybackslash}p{(\linewidth - 2\tabcolsep) * \real{0.5000}}
  >{\raggedright\arraybackslash}p{(\linewidth - 2\tabcolsep) * \real{0.5000}}@{}}
\toprule\noalign{}
\begin{minipage}[b]{\linewidth}\raggedright
Durée
\end{minipage} & \begin{minipage}[b]{\linewidth}\raggedright
Action avec l'enseignant
\end{minipage} \\
\midrule\noalign{}
\endhead
\bottomrule\noalign{}
\endlastfoot
10 min & Lire le rendu et votre question ; distinguer fourni, modifié,
exécuté et déclaré \\
15 min & Reformuler un scénario d'annulation et une sélection clavier \\
25 min & Raccorder l'interaction à une vraie zone du projet ; essayer
capture, relâchement, interruption et clavier \\
10 min & Adapter le geste à une carte et expliquer les données à
conserver \\
\end{longtable}
}

\subsection{Script enseignant --- séquence de 5
minutes}\label{script-enseignant-suxe9quence-de-5-minutes}

Les pauses et intonations ci-dessous sont proposées pour la lecture,
sans reproduction vérifiée d'une voix enregistrée.

\textbf{0:00--1:00 --- montrer.} « Je veux qu'on parte d'un geste
simple. Je prends l'objet, je le déplace et je relâche. {[}Pause ; faire
le geste.{]} Ce qui m'intéresse, c'est la donnée qui change. Dans ce
code, où se trouve sa position ? » Laisser l'élève pointer le champ ;
ouvrir ensuite \texttt{state.position}.

\textbf{1:00--2:00 --- prédire.} « Maintenant, je passe la souris sans
cliquer. {[}Accentuer : \emph{sans cliquer}.{]} Qu'est-ce qui devrait
changer ? On écrit notre idée avant de regarder. » Attendre une réponse,
puis montrer la condition sur \texttt{pointer\allowbreak{}Id} dans
\texttt{pointermove}.

\textbf{2:00--3:00 --- distinguer.} « Relâcher et être interrompu ont
une conséquence commune : le geste actif doit se terminer. La raison
reste différente. {[}Ralentir.{]} Une sortie de zone sous capture peut
continuer le geste. Regardons donc le nom de l'événement et
l'identifiant actif, ensemble. » Comparer deux états sans demander de
mémoriser toutes les API.

\textbf{3:00--4:00 --- essayer.} « Tu choisis le pas clavier. Prédis
deux déplacements, puis essaie. Je te laisse expliquer la différence. »
Laisser environ trente secondes à la manipulation et à la comparaison.

\textbf{4:00--5:00 --- transférer.} « Si cette sphère devient une carte
de ton site, quelle donnée gardes-tu ? Et si le geste s'interrompt,
comment éviter une carte qui reste accrochée ? » Accepter texte, dessin
ou code annoté. Si l'explication manque, revenir à un seul événement
plutôt que conclure à une maîtrise acquise.

\section{Repérer et travailler une
difficulté}\label{repuxe9rer-et-travailler-une-difficultuxe9}

La confusion à examiner est : \textbf{confondre survol, relâchement,
sortie de zone et annulation}. Demander un exemple précis, faire
retrouver une seule donnée, puis choisir une expérience plus petite si
nécessaire. Une réponse fluide de l'assistant peut devenir une hypothèse
de travail ; l'élève la relie au fichier avant de la conserver.

Retrouver le symbole \texttt{create\allowbreak{}Interaction} et appliquer le cadre
d'intervention commun à cette difficulté.

\section{Transfert et retour humain}\label{transfert-et-retour-humain}

Proposer \textbf{une carte déplaçable et sélectionnable au clavier}.
L'élève peut répondre par texte, schéma ou code annoté, puis essayer la
proposition. Relier le retour humain au scénario de transfert
effectivement examiné.

\chapter{Fiche enseignant 2 --- Organisation
Astro}\label{fiche-enseignant-2-organisation-astro}

\textbf{But de la revue :} L'élève retrouve trois responsabilités dans
les bons fichiers et distingue l'aperçu HTML de la compilation Astro.

\textbf{Préparation :} ouvrir le rendu sélectionné, le notebook et
\href{../../frontend/module-2/exploration-card.js}{exploration-card.js}.
Le \href{../../notebooks/instructor/module-2.ipynb}{corrigé séparé} est
une piste à examiner après la tentative. Retrouver la prédiction et la
version réellement choisie avant de comparer.

\section{Une activité avec l'élève et son
assistant}\label{une-activituxe9-avec-luxe9luxe8ve-et-son-assistant-1}

L'élève apporte trois contenus personnels. Avec l'assistant, il indique
quel fichier fournit la section et lequel remplit les boutons. Il
modifie le préfixe puis explique la partie conservée.

L'assistant peut expliquer ou fournir un exemple complet. L'élève
choisit ensuite sa modification, nomme l'aide reçue et montre un
résultat. La question à poser est : \textbf{« Où changes-tu le titre,
les données et la réaction au clic ? »}

Comparer \textbf{descriptionPrefix : Mon observation → Ma décision
expliquée}, en conservant \textbf{identifiants, nombre de boutons et
sélection}. Le rendu sélectionné peut prendre la forme de arborescence
annotée, deux fichiers exportés et choix de contenu. La limite à garder
visible est : \textbf{l'aperçu Colab ne compile pas
ExplorationCard.astro}.

\section{Déroulé et script de
séance}\label{duxe9rouluxe9-et-script-de-suxe9ance-1}

{\def\LTcaptype{none} % do not increment counter
\begin{longtable}[]{@{}
  >{\raggedright\arraybackslash}p{(\linewidth - 2\tabcolsep) * \real{0.5000}}
  >{\raggedright\arraybackslash}p{(\linewidth - 2\tabcolsep) * \real{0.5000}}@{}}
\toprule\noalign{}
\begin{minipage}[b]{\linewidth}\raggedright
Durée
\end{minipage} & \begin{minipage}[b]{\linewidth}\raggedright
Action avec l'enseignant
\end{minipage} \\
\midrule\noalign{}
\endhead
\bottomrule\noalign{}
\endlastfoot
10 min & Examiner le rendu et l'arborescence proposée \\
15 min & Clarifier contenu, structure et comportement sur un exemple
personnel \\
25 min & Raccorder la carte dans le projet Astro ; vérifier titre,
sélection, clavier et compilation \\
10 min & Réutiliser la carte pour un nouveau sujet sans recopier la
logique \\
\end{longtable}
}

\subsection{Script enseignant --- séquence de 5
minutes}\label{script-enseignant-suxe9quence-de-5-minutes-1}

Les intonations sont des propositions de lecture.

\textbf{0:00--1:00 --- partir de la carte.} « On a une carte, un titre
et trois choix. Je veux changer le sujet. {[}Pause.{]} Avant de toucher
au code, quel fichier ouvre-t-on ? Montre-moi ton raisonnement avec le
titre que tu veux remplacer. » Laisser l'élève choisir un fichier et
expliquer pourquoi.

\textbf{1:00--2:00 --- suivre une donnée.} « Ici, Astro fournit une
place : une section, une liste, une description. Le contrôleur remplit
cette place. {[}Accentuer : \emph{la même donnée, un rôle précis}.{]}
Suivons un identifiant jusqu'au bouton. » Montrer les données,
\texttt{dataset.learning\allowbreak{}Id} et le rappel \texttt{on\allowbreak{}Select}.

\textbf{2:00--3:00 --- comparer.} « Tu as changé un préfixe. Qu'as-tu vu
? Et qu'est-ce qui aurait dû rester identique ? » Attendre le constat de
l'élève. En cas d'écart, regarder le fichier exécuté plutôt que supposer
une erreur de compréhension.

\textbf{3:00--4:00 --- poser la limite.} « Le carnet nous donne un
aperçu HTML. Pour savoir si le vrai composant Astro se construit, on
doit examiner sa compilation. {[}Ralentir.{]} Ce sont deux vérifications
utiles, avec deux objets différents. » Faire nommer l'objet vérifié par
chaque étape.

\textbf{4:00--5:00 --- transférer.} « Imagine une carte pour trois
livres ou trois activités. Quelles données changes-tu ? Quelle logique
peux-tu garder ? » Laisser annoter trois modifications puis expliquer la
partie conservée. La répétition d'une phrase du cours ne suffit pas :
l'élève doit montrer son nouveau contenu.

\section{Repérer et travailler une
difficulté}\label{repuxe9rer-et-travailler-une-difficultuxe9-1}

La confusion à examiner est : \textbf{attribuer au composant Astro un
comportement produit par le contrôleur navigateur}. Demander un exemple
précis, faire retrouver une seule donnée, puis choisir une expérience
plus petite si nécessaire. Une réponse fluide de l'assistant peut
devenir une hypothèse de travail ; l'élève la relie au fichier avant de
la conserver.

Retrouver le symbole \texttt{create\allowbreak{}Cards} et appliquer le cadre
d'intervention commun à cette difficulté.

\section{Transfert et retour humain}\label{transfert-et-retour-humain-1}

Proposer \textbf{une carte présentant des livres, des étapes ou des
activités}. L'élève peut répondre par texte, schéma ou code annoté, puis
essayer la proposition. Relier le retour humain au scénario de transfert
effectivement examiné.

\chapter{Fiche enseignant 3 --- Scène Three.js et
coordonnées}\label{fiche-enseignant-3-scuxe8ne-three.js-et-coordonnuxe9es}

\textbf{But de la revue :} L'élève nomme la caméra et retrouve la
position conservée ; il utilise aussi le contenu HTML.

\textbf{Préparation :} ouvrir le rendu sélectionné, le notebook et
\href{../../frontend/module-3/scene-controller.js}{scene-controller.js}.
Le \href{../../notebooks/instructor/module-3.ipynb}{corrigé séparé} est
une piste à examiner après la tentative. Retrouver la prédiction et la
version réellement choisie avant de comparer.

\section{Une activité avec l'élève et son
assistant}\label{une-activituxe9-avec-luxe9luxe8ve-et-son-assistant-2}

Montrer le même élément dans la scène et la liste. L'élève prédit
l'effet du recul de caméra, sélectionne par les deux représentations et
retrouve l'identifiant commun.

L'assistant peut expliquer ou fournir un exemple complet. L'élève
choisit ensuite sa modification, nomme l'aide reçue et montre un
résultat. La question à poser est : \textbf{« L'objet semble plus petit
: quelle donnée a réellement changé ? »}

Comparer \textbf{cameraDistance : 5 → 7}, en conservant \textbf{objets,
identifiants et positions}. Le rendu sélectionné peut prendre la forme
de deux paramètres de caméra, sélection, disponibilité WebGL et
explication. La limite à garder visible est : \textbf{une scène
indisponible donne une observation de repli HTML, pas une observation
3D}.

\section{Déroulé et script de
séance}\label{duxe9rouluxe9-et-script-de-suxe9ance-2}

{\def\LTcaptype{none} % do not increment counter
\begin{longtable}[]{@{}
  >{\raggedright\arraybackslash}p{(\linewidth - 2\tabcolsep) * \real{0.5000}}
  >{\raggedright\arraybackslash}p{(\linewidth - 2\tabcolsep) * \real{0.5000}}@{}}
\toprule\noalign{}
\begin{minipage}[b]{\linewidth}\raggedright
Durée
\end{minipage} & \begin{minipage}[b]{\linewidth}\raggedright
Action avec l'enseignant
\end{minipage} \\
\midrule\noalign{}
\endhead
\bottomrule\noalign{}
\endlastfoot
10 min & Examiner la comparaison et le statut WebGL \\
15 min & Reformuler la distinction objet/caméra/projection \\
25 min & Raccorder sélection 3D et fiche HTML ; essayer
redimensionnement et parcours accessible \\
10 min & Transférer les mêmes identifiants à une autre représentation \\
\end{longtable}
}

\subsection{Script enseignant --- séquence de 5
minutes}\label{script-enseignant-suxe9quence-de-5-minutes-2}

Les pauses et intonations proposées servent à laisser une place à
l'observation.

\textbf{0:00--1:00 --- observer.} « Je veux qu'on regarde le même objet
sous deux vues. {[}Pause ; montrer la sphère et son bouton.{]} Quelle
information nous dit qu'il s'agit du même élément ? Montre-moi cet
identifiant dans les données. » Laisser chercher avant d'ouvrir le
rappel de sélection.

\textbf{1:00--2:00 --- prédire.} « Je vais reculer la caméra.
{[}Accentuer : \emph{la caméra}.{]} Est-ce que je dois changer la
position enregistrée de l'objet ? Écris ta prédiction. » Attendre la
distinction entre données et image.

\textbf{2:00--3:00 --- comparer.} « Passons de 5 à 7. Décris ce que tu
vois avant de l'expliquer. Ensuite, vérifions les valeurs conservées. »
Donner du temps à la comparaison. Une observation absente devient une
limite, pas un résultat supposé.

\textbf{3:00--4:00 --- retrouver l'information.} « Si ton appareil ne
dessine pas cette scène, comment gardes-tu accès au même contenu ?
{[}Ralentir.{]} Le bouton HTML représente un vrai parcours de sélection.
Essaie-le et explique ce qu'il garde. » Vérifier l'identifiant et la
description sans réduire l'exercice à l'effet visuel.

\textbf{4:00--5:00 --- transférer.} « Pour ton projet, pourrais-tu
présenter trois étapes dans une scène et dans une liste ? Quelle donnée
doit rester commune ? » Laisser dessiner un exemple. Demander une
justification de l'intérêt de la 3D et conserver la question si cet
intérêt reste incertain.

\section{Repérer et travailler une
difficulté}\label{repuxe9rer-et-travailler-une-difficultuxe9-2}

La confusion à examiner est : \textbf{confondre position enregistrée,
position dans la scène et taille projetée}. Demander un exemple précis,
faire retrouver une seule donnée, puis choisir une expérience plus
petite si nécessaire. Une réponse fluide de l'assistant peut devenir une
hypothèse de travail ; l'élève la relie au fichier avant de la
conserver.

Retrouver le symbole \texttt{create\allowbreak{}Scene\allowbreak{}Controller} et appliquer le
cadre d'intervention commun à cette difficulté.

\section{Transfert et retour humain}\label{transfert-et-retour-humain-2}

Proposer \textbf{trois étapes consultables dans une scène et une liste}.
L'élève peut répondre par texte, schéma ou code annoté, puis essayer la
proposition. Relier le retour humain au scénario de transfert
effectivement examiné.

\chapter{Fiche enseignant 4 --- Vitesse et
inertie}\label{fiche-enseignant-4-vitesse-et-inertie}

\textbf{But de la revue :} L'élève définit les unités, décrit le rôle de
vx/vy, distingue déplacement clavier sans élan et lancer au pointeur, et
reconnaît la différence entre deux gestes humains et deux entrées
contrôlées.

\textbf{Préparation :} ouvrir le rendu sélectionné, le notebook et
\href{../../frontend/module-4/inertia.js}{inertia.js}. Le
\href{../../notebooks/instructor/module-4.ipynb}{corrigé séparé} est une
piste à examiner après la tentative. Retrouver la prédiction et la
version réellement choisie avant de comparer.

\section{Une activité avec l'élève et son
assistant}\label{une-activituxe9-avec-luxe9luxe8ve-et-son-assistant-3}

L'élève calcule un déplacement simple avec dt, prédit l'effet d'un
amortissement supérieur et compare deux lancers au pointeur. Il utilise
aussi les flèches pour déplacer de 0,04 avec une vitesse nulle, puis
explique pourquoi leur relâchement ne lance pas l'objet. Espace
fige/reprend; Échap annule le geste. Le pointeur passif ne réagrippe pas
l'objet lancé.

L'assistant peut expliquer ou fournir un exemple complet. L'élève
choisit ensuite sa modification, nomme l'aide reçue et montre un
résultat. La question à poser est : \textbf{« Après le relâchement du
pointeur, quelle donnée conserve le mouvement ? »}

Comparer \textbf{damping : 0.65 → 1.3}, en conservant
\textbf{restitution 0.9, zone et geste aussi comparable que possible}.
Le rendu sélectionné peut prendre la forme de paramètres, trajectoires
déclarées, unités et conditions du lancer. La limite à garder visible
est : \textbf{les gestes humains ne garantissent pas une vitesse
initiale identique; les flèches ne constituent pas une comparaison
d'inertie}.

\section{Déroulé et script de
séance}\label{duxe9rouluxe9-et-script-de-suxe9ance-3}

{\def\LTcaptype{none} % do not increment counter
\begin{longtable}[]{@{}
  >{\raggedright\arraybackslash}p{(\linewidth - 2\tabcolsep) * \real{0.5000}}
  >{\raggedright\arraybackslash}p{(\linewidth - 2\tabcolsep) * \real{0.5000}}@{}}
\toprule\noalign{}
\begin{minipage}[b]{\linewidth}\raggedright
Durée
\end{minipage} & \begin{minipage}[b]{\linewidth}\raggedright
Action avec l'enseignant
\end{minipage} \\
\midrule\noalign{}
\endhead
\bottomrule\noalign{}
\endlastfoot
10 min & Examiner les deux essais et la comparabilité des gestes \\
15 min & Reformuler position, vitesse et \texttt{dt} avec un exemple
chiffré \\
25 min & Raccorder lancer et rebonds dans le navigateur cible ; vérifier
clavier, annulation et mode figé \\
10 min & Appliquer le même mouvement à une carte HTML ou une caméra \\
\end{longtable}
}

\subsection{Script enseignant --- séquence de 5
minutes}\label{script-enseignant-suxe9quence-de-5-minutes-3}

Les indications de lecture sont proposées.

\textbf{0:00--1:00 --- poser la question.} « Je relâche, et l'objet
continue. {[}Pause ; lancer.{]} Mon pointeur n'est plus actif : quelle
donnée reste dans l'objet ? Montre-moi cette donnée. » Attendre
\texttt{vx/\allowbreak{}vy} ou une explication équivalente avant de nommer la
vitesse.

\textbf{1:00--2:00 --- définir les unités.} « Imaginons 0,2 unité par
seconde pendant 0,05 seconde. {[}Ralentir.{]} On avance de 0,01 unité.
Le temps compte. Avec une autre durée, le déplacement change, même si la
vitesse est identique. » Faire reformuler le calcul avec un autre
nombre.

\textbf{2:00--3:00 --- prédire et essayer.} « Je garde la restitution et
je change un amortissement. Avant de lancer, lequel devrait ralentir
plus vite ? » Écrire la prédiction ; effectuer la comparaison. Laisser
l'élève décrire avant d'expliquer.

\textbf{3:00--4:00 --- examiner la limite.} « Mes gestes ne sont pas
exactement les mêmes. Alors, quel résultat ai-je observé et quelle
partie faudrait-il contrôler davantage ? » Montrer la vitesse initiale
si elle est disponible. Conserver un résultat incertain comme tel.

\textbf{4:00--5:00 --- transférer.} « Maintenant, tu veux une carte qui
glisse après le geste. Quelle partie gardes-tu sans Three.js ? Et
pourquoi le simple passage de la souris ne doit-il pas reprendre la
carte ? » Laisser relier geste actif et mouvement autonome, puis
préparer un essai dans un nouveau contexte.

\section{Repérer et travailler une
difficulté}\label{repuxe9rer-et-travailler-une-difficultuxe9-3}

La confusion à examiner est : \textbf{traiter une position, une vitesse
et une durée comme la même grandeur}. Demander un exemple précis, faire
retrouver une seule donnée, puis choisir une expérience plus petite si
nécessaire. Une réponse fluide de l'assistant peut devenir une hypothèse
de travail ; l'élève la relie au fichier avant de la conserver.

Retrouver le symbole \texttt{create\allowbreak{}Inertia} et appliquer le cadre
d'intervention commun à cette difficulté.

\section{Transfert et retour humain}\label{transfert-et-retour-humain-3}

Proposer \textbf{une carte qui glisse après relâchement}. L'élève peut
répondre par texte, schéma ou code annoté, puis essayer la proposition.
Relier le retour humain au scénario de transfert effectivement examiné.

\chapter{Fiche enseignant 5 --- Interpolation et
élasticité}\label{fiche-enseignant-5-interpolation-et-uxe9lasticituxe9}

\textbf{But de la revue :} L'élève distingue valeur présente, cible et
vitesse ; il montre le gel sans disparition de l'image ni rattrapage du
temps suspendu.

\textbf{Préparation :} ouvrir le rendu sélectionné, le notebook et
\href{../../frontend/module-5/transitions.js}{transitions.js}. Le
\href{../../notebooks/instructor/module-5.ipynb}{corrigé séparé} est une
piste à examiner après la tentative. Retrouver la prédiction et la
version réellement choisie avant de comparer.

\section{Une activité avec l'élève et son
assistant}\label{une-activituxe9-avec-luxe9luxe8ve-et-son-assistant-4}

Comparer une approche de cible et un dépassement. L'élève prédit la
nouvelle trajectoire, fige l'image en cours de mouvement, puis explique
valeur, cible et vitesse lors de la reprise.

L'assistant peut expliquer ou fournir un exemple complet. L'élève
choisit ensuite sa modification, nomme l'aide reçue et montre un
résultat. La question à poser est : \textbf{« Que peut-il rester en
mouvement lorsque la valeur arrive à sa cible ? »}

Comparer \textbf{damping : 12 → 5}, en conservant \textbf{stiffness 36
et cible}. Le rendu sélectionné peut prendre la forme de paramètre
choisi, observation d'oscillation et scénario de gel/reprise. La limite
à garder visible est : \textbf{un essai réussi ne qualifie pas toutes
les valeurs possibles de ressort}.

\section{Déroulé et script de
séance}\label{duxe9rouluxe9-et-script-de-suxe9ance-4}

{\def\LTcaptype{none} % do not increment counter
\begin{longtable}[]{@{}
  >{\raggedright\arraybackslash}p{(\linewidth - 2\tabcolsep) * \real{0.5000}}
  >{\raggedright\arraybackslash}p{(\linewidth - 2\tabcolsep) * \real{0.5000}}@{}}
\toprule\noalign{}
\begin{minipage}[b]{\linewidth}\raggedright
Durée
\end{minipage} & \begin{minipage}[b]{\linewidth}\raggedright
Action avec l'enseignant
\end{minipage} \\
\midrule\noalign{}
\endhead
\bottomrule\noalign{}
\endlastfoot
10 min & Examiner la comparaison et les valeurs conservées \\
15 min & Reformuler cible, vitesse, dépassement et stabilisation \\
25 min & Raccorder une transition du projet ; tester gel, reprise,
clavier et nettoyage \\
10 min & Transférer le mécanisme à une fiche ou une caméra \\
\end{longtable}
}

\subsection{Script enseignant --- séquence de 5
minutes}\label{script-enseignant-suxe9quence-de-5-minutes-4}

Les intonations sont proposées pour laisser du temps à la prédiction et
à l'essai.

\textbf{0:00--1:00 --- montrer.} « Je veux qu'on distingue deux retours
vers une cible. {[}Pause ; montrer deux trajectoires.{]} Les deux
reviennent, mais l'une dépasse. Quelle donnée peut continuer à évoluer
lorsque la valeur arrive à sa cible ? » Laisser l'élève formuler une
idée avant de montrer la vitesse.

\textbf{1:00--2:00 --- relier au code.} « Dans ce code, la distance à la
cible attire ; la vitesse conservée peut nous faire passer de l'autre
côté. {[}Ralentir.{]} L'amortissement freine cette vitesse. Suivons une
seule ligne plutôt que toute l'animation. » Montrer les deux termes de
l'accélération.

\textbf{2:00--3:00 --- prédire.} « On garde 36 pour la raideur et on
passe de 12 à 5 pour l'amortissement. Quel changement attends-tu ? »
Écrire la réponse puis lancer. Donner environ trente secondes à
l'observation.

\textbf{3:00--4:00 --- figer.} « Maintenant, je fige l'image.
{[}Pause.{]} Où doit rester la valeur ? On peut vérifier cette valeur
sans tout enlever. À la reprise, le temps suspendu ne doit pas devenir
un saut. » Faire expliquer les données conservées et la chronologie.

\textbf{4:00--5:00 --- transférer.} « Dans ton projet, quelle transition
profiterait de ce mécanisme ? Et à quel moment préfères-tu un changement
immédiat ? » Laisser proposer une fiche ou une caméra, avec une
justification d'usage. Une préférence esthétique et une vérification
technique sont deux traces distinctes.

\section{Repérer et travailler une
difficulté}\label{repuxe9rer-et-travailler-une-difficultuxe9-4}

La confusion à examiner est : \textbf{identifier un mécanisme uniquement
à l'apparence douce de son animation}. Demander un exemple précis, faire
retrouver une seule donnée, puis choisir une expérience plus petite si
nécessaire. Une réponse fluide de l'assistant peut devenir une hypothèse
de travail ; l'élève la relie au fichier avant de la conserver.

Retrouver le symbole \texttt{create\allowbreak{}Transition} et appliquer le cadre
d'intervention commun à cette difficulté.

\section{Transfert et retour humain}\label{transfert-et-retour-humain-4}

Proposer \textbf{l'ouverture d'une fiche ou le déplacement d'une
caméra}. L'élève peut répondre par texte, schéma ou code annoté, puis
essayer la proposition. Relier le retour humain au scénario de transfert
effectivement examiné.

\chapter{Fiche enseignant 6 --- Particules et
ressources}\label{fiche-enseignant-6-particules-et-ressources}

\textbf{But de la revue :} L'élève sépare population et performance,
explique la réutilisation d'un emplacement et montre le nettoyage.

\textbf{Préparation :} ouvrir le rendu sélectionné, le notebook et
\href{../../frontend/module-6/particle-layer.js}{particle-layer.js}. Le
\href{../../notebooks/instructor/module-6.ipynb}{corrigé séparé} est une
piste à examiner après la tentative. Retrouver la prédiction et la
version réellement choisie avant de comparer.

\section{Une activité avec l'élève et son
assistant}\label{une-activituxe9-avec-luxe9luxe8ve-et-son-assistant-5}

L'élève annote une particule et prédit ce qui augmente avec la
population. Il compare le même scénario, indique la source d'une mesure
disponible et conserve un libellé avec l'effet désactivé.

L'assistant peut expliquer ou fournir un exemple complet. L'élève
choisit ensuite sa modification, nomme l'aide reçue et montre un
résultat. La question à poser est : \textbf{« Que mesure ton
chronomètre, et quel résultat reste à mesurer dans le navigateur cible ?
»}

Comparer \textbf{count : 64 → 128, puis restauration à 64}, en
conservant \textbf{lifetime 3 et seed 17}. Le rendu sélectionné peut
prendre la forme de population, durée, graine, appareil, mesure
disponible et limite. La limite à garder visible est : \textbf{le temps
Python ou l'aperçu Canvas ne mesure pas directement un rendu GPU
Three.js}.

\section{Déroulé et script de
séance}\label{duxe9rouluxe9-et-script-de-suxe9ance-5}

{\def\LTcaptype{none} % do not increment counter
\begin{longtable}[]{@{}
  >{\raggedright\arraybackslash}p{(\linewidth - 2\tabcolsep) * \real{0.5000}}
  >{\raggedright\arraybackslash}p{(\linewidth - 2\tabcolsep) * \real{0.5000}}@{}}
\toprule\noalign{}
\begin{minipage}[b]{\linewidth}\raggedright
Durée
\end{minipage} & \begin{minipage}[b]{\linewidth}\raggedright
Action avec l'enseignant
\end{minipage} \\
\midrule\noalign{}
\endhead
\bottomrule\noalign{}
\endlastfoot
10 min & Examiner paramètres, appareil et source de mesure \\
15 min & Clarifier réutilisation, durée de vie et nettoyage \\
25 min & Mesurer dans le navigateur cible ; tester une couche
facultative et sa fermeture \\
10 min & Transférer l'effet à une sélection, avec libellé et mode
désactivé \\
\end{longtable}
}

\subsection{Script enseignant --- séquence de 5
minutes}\label{script-enseignant-suxe9quence-de-5-minutes-5}

Les pauses et intonations ci-dessous sont proposées.

\textbf{0:00--1:00 --- voir les données.} « Je veux qu'on regarde un
point comme une fiche de données. Il a une position, une vitesse et un
âge. {[}Pause.{]} Si j'en affiche 64, où est-ce que je retrouve ce
nombre dans le code ? » Laisser chercher \texttt{count} et la création
de la liste.

\textbf{1:00--2:00 --- prédire.} « On va en mettre 128 avec la même
durée de vie. Qu'est-ce qui augmente certainement ? Et quelle
performance reste à mesurer ? » Accepter une prédiction nuancée ;
distinguer population et durée de rendu.

\textbf{2:00--3:00 --- essayer.} « Exécutons le même scénario. Décris
d'abord ce que tu vois. {[}Pause d'observation.{]} Maintenant, quelle
mesure peux-tu réellement citer ? » Si aucune métrique n'est disponible,
consigner ce fait sans inventer une cadence.

\textbf{3:00--4:00 --- fermer.} « Ces éléments peuvent finir leur vie
sans créer une liste infinie. Regardons ce qui est réutilisé. Puis
fermons la couche : qui libère ses données, qui libère le matériel
graphique ? » Relier \texttt{dispose()} aux responsabilités séparées.

\textbf{4:00--5:00 --- choisir l'usage.} « Pour ton projet, je te
propose un effet qui souligne une sélection. Comment garder le texte
lisible quand l'effet est désactivé ? » Laisser dessiner le parcours
statique. Le choix d'un effet devient une décision expliquée, pas une
obligation de multiplier les points.

\section{Repérer et travailler une
difficulté}\label{repuxe9rer-et-travailler-une-difficultuxe9-5}

La confusion à examiner est : \textbf{confondre nombre de données,
nombre de maillages et coût GPU}. Demander un exemple précis, faire
retrouver une seule donnée, puis choisir une expérience plus petite si
nécessaire. Une réponse fluide de l'assistant peut devenir une hypothèse
de travail ; l'élève la relie au fichier avant de la conserver.

Retrouver le symbole \texttt{create\allowbreak{}Particle\allowbreak{}Layer} et appliquer le cadre
d'intervention commun à cette difficulté.

\section{Transfert et retour humain}\label{transfert-et-retour-humain-5}

Proposer \textbf{un effet facultatif indiquant une sélection}. L'élève
peut répondre par texte, schéma ou code annoté, puis essayer la
proposition. Relier le retour humain au scénario de transfert
effectivement examiné.

\chapter{Fiche enseignant 7 --- Parcours, états et
mémoire}\label{fiche-enseignant-7-parcours-uxe9tats-et-muxe9moire}

\textbf{But de la revue :} L'élève distingue historique interne et
navigation navigateur, conserve vue et sélection et explique le gel de
elapsed.

\textbf{Préparation :} ouvrir le rendu sélectionné, le notebook et
\href{../../frontend/module-7/interaction-flow.js}{interaction-flow.js}.
Le \href{../../notebooks/instructor/module-7.ipynb}{corrigé séparé} est
une piste à examiner après la tentative. Retrouver la prédiction et la
version réellement choisie avant de comparer.

\section{Une activité avec l'élève et son
assistant}\label{une-activituxe9-avec-luxe9luxe8ve-et-son-assistant-6}

L'élève trace cinq transitions pour un budget de quatre, retrouve
l'entrée retirée et vérifie la sélection. Dans le vrai projet, il fait
un aller-retour navigateur et explique l'URL restaurée.

L'assistant peut expliquer ou fournir un exemple complet. L'élève
choisit ensuite sa modification, nomme l'aide reçue et montre un
résultat. La question à poser est : \textbf{« Au retour de la fiche,
quelles données doivent être retrouvées ensemble ? »}

Comparer \textbf{historyLimit : 32 → 4}, en conservant \textbf{séquence
précise de cinq changements}. Le rendu sélectionné peut prendre la forme
de état initial, séquence, entrée retirée et état restauré. La limite à
garder visible est : \textbf{host: null dans Colab ne valide pas le
bouton Retour de la page Colab}.

\section{Déroulé et script de
séance}\label{duxe9rouluxe9-et-script-de-suxe9ance-6}

{\def\LTcaptype{none} % do not increment counter
\begin{longtable}[]{@{}
  >{\raggedright\arraybackslash}p{(\linewidth - 2\tabcolsep) * \real{0.5000}}
  >{\raggedright\arraybackslash}p{(\linewidth - 2\tabcolsep) * \real{0.5000}}@{}}
\toprule\noalign{}
\begin{minipage}[b]{\linewidth}\raggedright
Durée
\end{minipage} & \begin{minipage}[b]{\linewidth}\raggedright
Action avec l'enseignant
\end{minipage} \\
\midrule\noalign{}
\endhead
\bottomrule\noalign{}
\endlastfoot
10 min & Examiner la séquence et la question de restauration \\
15 min & Clarifier geste, parcours, historique interne et URL \\
25 min & Tester fiche, retour navigateur, interruption et gel dans le
projet \\
10 min & Transférer le parcours à une galerie ou un explorateur \\
\end{longtable}
}

\subsection{Script enseignant --- séquence de 5
minutes}\label{script-enseignant-suxe9quence-de-5-minutes-6}

Les intonations sont proposées pour accompagner la démonstration.

\textbf{0:00--1:00 --- formuler l'attente.} « Je sélectionne le deuxième
élément, puis j'ouvre sa fiche. {[}Pause.{]} Quand je reviens, où
dois-je me retrouver ? Je veux qu'on nomme la vue et la sélection avant
de coder. » Laisser tracer les états.

\textbf{1:00--2:00 --- distinguer.} « Le pointeur peut être relâché,
alors que la fiche reste ouverte. {[}Accentuer : \emph{deux états
différents}.{]} Quelle donnée appartient au geste ? Quelle donnée
appartient au parcours ? » Revenir aux fichiers des Modules 1 et 7 si
les rôles sont mélangés.

\textbf{2:00--3:00 --- suivre les entrées.} « On garde quatre entrées et
on fait cinq changements. Montre-moi celle qui sort. » Laisser manipuler
et expliquer. Demander aussi ce que conserve la sélection courante,
indépendamment du budget de diagnostic.

\textbf{3:00--4:00 --- vérifier le vrai navigateur.} « Le carnet nous
montre un historique interne. Ici, dans le projet, l'URL et le bouton
Retour doivent restaurer l'état. {[}Ralentir.{]} On vérifie donc ce
parcours dans son navigateur cible. » Faire un aller-retour et consigner
un éventuel écart.

\textbf{4:00--5:00 --- transférer.} « Si cette vue devient une galerie,
que faut-il conserver lorsque tu ouvres une image puis reviens ? Et si
la page est figée, quelle horloge doit rester en place ? » Accepter
schéma ou code annoté, avec un essai nouveau avant une conclusion de
compréhension.

\section{Repérer et travailler une
difficulté}\label{repuxe9rer-et-travailler-une-difficultuxe9-6}

La confusion à examiner est : \textbf{mélanger l'état du geste, l'état
de la vue et l'historique global du navigateur}. Demander un exemple
précis, faire retrouver une seule donnée, puis choisir une expérience
plus petite si nécessaire. Une réponse fluide de l'assistant peut
devenir une hypothèse de travail ; l'élève la relie au fichier avant de
la conserver.

Retrouver le symbole \texttt{create\allowbreak{}Interaction\allowbreak{}Flow} et appliquer le
cadre d'intervention commun à cette difficulté.

\section{Transfert et retour humain}\label{transfert-et-retour-humain-6}

Proposer \textbf{une galerie qui revient à la bonne sélection}. L'élève
peut répondre par texte, schéma ou code annoté, puis essayer la
proposition. Relier le retour humain au scénario de transfert
effectivement examiné.

\chapter{Fiche enseignant 8 --- Assistant, preuves et
capacités}\label{fiche-enseignant-8-assistant-preuves-et-capacituxe9s}

\textbf{But de la revue :} L'élève attribue correctement les preuves,
borne la lecture et explique sa contribution dans le projet assemblé.

\textbf{Préparation :} ouvrir le rendu sélectionné, le notebook et
\href{../../frontend/module-8/register-capabilities.js}{register-capabilities.js}.
Le \href{../../notebooks/instructor/module-8.ipynb}{corrigé séparé} est
une piste à examiner après la tentative. Retrouver la prédiction et la
version réellement choisie avant de comparer.

\section{Une activité avec l'élève et son
assistant}\label{une-activituxe9-avec-luxe9luxe8ve-et-son-assistant-7}

L'élève examine la trace préparée, montre l'objection liée à la
permission ou à la révision et définit son contenu partageable. Pendant
la séance, il retrouve une vraie modification dans chacun des huit
exports assemblés.

L'assistant peut expliquer ou fournir un exemple complet. L'élève
choisit ensuite sa modification, nomme l'aide reçue et montre un
résultat. La question à poser est : \textbf{« Quelle preuve distingue
une trace préparée, un appel réel et une revue humaine ? »}

Comparer \textbf{borne de snapshot : 12000 → 6000 caractères JSON}, en
conservant \textbf{permission, révision, revérification et nettoyage}.
Le rendu sélectionné peut prendre la forme de périmètre de capacité,
objection, code choisi et lien vers les exports. La limite à garder
visible est : \textbf{le notebook ne réalise pas l'appel natif ;
student\_get\_learning\_snapshot diffère du registre Orbit de vingt-cinq
outils}.

\section{Déroulé et script de
séance}\label{duxe9rouluxe9-et-script-de-suxe9ance-7}

{\def\LTcaptype{none} % do not increment counter
\begin{longtable}[]{@{}
  >{\raggedright\arraybackslash}p{(\linewidth - 2\tabcolsep) * \real{0.5000}}
  >{\raggedright\arraybackslash}p{(\linewidth - 2\tabcolsep) * \real{0.5000}}@{}}
\toprule\noalign{}
\begin{minipage}[b]{\linewidth}\raggedright
Durée
\end{minipage} & \begin{minipage}[b]{\linewidth}\raggedright
Action avec l'enseignant
\end{minipage} \\
\midrule\noalign{}
\endhead
\bottomrule\noalign{}
\endlastfoot
10 min & Examiner la trace, le code exporté et les permissions
prévues \\
15 min & Clarifier citation, portée, révision et autorité \\
25 min & Utiliser les vrais outils d'Orbit, puis raccorder vos huit
productions \\
10 min & Transférer une capacité bornée à votre projet et expliquer le
socle obtenu \\
\end{longtable}
}

\subsection{La découverte de fin de
leçon}\label{la-duxe9couverte-de-fin-de-leuxe7on}

Vous avez conservé huit briques. \textbf{Elles peuvent maintenant former
un socle de frontend :} geste et état, carte, scène, inertie,
transition, particules, parcours et capacité.

Pendant les 25 minutes de construction du webinaire du Module 8, dans la
partie finale du notebook 8, charger vos \textbf{huit ZIP réels}. Le
raccordement fourni vérifie les manifestes et prépare un projet avec vos
fichiers. Il ne comble pas une absence avec un corrigé caché. Si un
module manque ou si une empreinte ne correspond pas, conserver l'erreur
et réexporter le bon rendu.

Le code d'intégration est fourni et attribué au cours dans
\href{../../assembly/README.md}{assembly/}. Vos modifications sont dans
les dossiers \texttt{src/\allowbreak{}learning/\allowbreak{}module-N/\allowbreak{}}. L'assemblage du ZIP ne
compile pas Astro : la compilation et le parcours du projet sont
vérifiés séparément. Une modification hors du contrat appelle une
adaptation expliquée ; elle ne justifie pas une substitution silencieuse
du travail.

\subsection{Script enseignant --- séquence de 5
minutes}\label{script-enseignant-suxe9quence-de-5-minutes-7}

Les intonations sont proposées ; elles ne constituent pas un clonage de
voix.

\textbf{0:00--1:00 --- choisir ce qui est partagé.} « Je veux que ton
assistant puisse t'aider sur un contenu que tu choisis. {[}Pause.{]}
Quelle fiche lui donnes-tu, et quelle révision ? Montrons ces deux
conditions dans le code. » Laisser l'élève borner le périmètre avant
l'appel.

\textbf{1:00--2:00 --- lire la trace.} « Voici une réponse préparée qui
dit READY. {[}Accentuer : \emph{préparée}.{]} Ce mot peut-il prouver que
le navigateur a réellement exécuté l'outil ? » Attendre la distinction
puis examiner les champs incohérents avec la permission ou la révision.

\textbf{2:00--3:00 --- vérifier l'autorité.} « Un agent peut expliquer,
calculer et proposer. Une revue humaine demande une décision réellement
humaine. {[}Ralentir.{]} Regardons aussi le passage et les conditions
avant d'accepter la conclusion. » Demander une objection précise plutôt
qu'un sentiment général de confiance.

\textbf{3:00--4:00 --- découvrir le raccordement.} « Tu as gardé tes
huit productions. Maintenant, on va les raccorder. Ce sont tes exports
qui entrent dans le projet ; le code d'assemblage est fourni. » Montrer
un chemin réellement importé et laisser l'élève retrouver une
modification.

\textbf{4:00--5:00 --- choisir la suite.} « Ce socle appartient
maintenant à ton projet. Quel usage veux-tu en faire ? Et quelle
capacité utile donneras-tu à ton assistant, avec un partage précis ? »
Faire expliquer un nouvel usage et un contrôle. Garder les incertitudes
pour la revue finale.

\section{Repérer et travailler une
difficulté}\label{repuxe9rer-et-travailler-une-difficultuxe9-7}

La confusion à examiner est : \textbf{prendre READY, une empreinte ou
une proposition pour une approbation humaine}. Demander un exemple
précis, faire retrouver une seule donnée, puis choisir une expérience
plus petite si nécessaire. Une réponse fluide de l'assistant peut
devenir une hypothèse de travail ; l'élève la relie au fichier avant de
la conserver.

Retrouver le symbole \texttt{register\allowbreak{}Learning\allowbreak{}Capability} et appliquer le
cadre d'intervention commun à cette difficulté.

\section{Transfert et retour humain}\label{transfert-et-retour-humain-7}

Proposer \textbf{une capacité bornée utile à son propre projet}. L'élève
peut répondre par texte, schéma ou code annoté, puis essayer la
proposition. Relier le retour humain au scénario de transfert
effectivement examiné.

\chapter{Fiche enseignant projet --- cadrage et
clôture}\label{fiche-enseignant-projet-cadrage-et-cluxf4ture}

\section{Une heure de cadrage}\label{une-heure-de-cadrage}

{\def\LTcaptype{none} % do not increment counter
\begin{longtable}[]{@{}
  >{\raggedleft\arraybackslash}p{(\linewidth - 4\tabcolsep) * \real{0.4000}}
  >{\raggedright\arraybackslash}p{(\linewidth - 4\tabcolsep) * \real{0.3000}}
  >{\raggedright\arraybackslash}p{(\linewidth - 4\tabcolsep) * \real{0.3000}}@{}}
\toprule\noalign{}
\begin{minipage}[b]{\linewidth}\raggedleft
Durée
\end{minipage} & \begin{minipage}[b]{\linewidth}\raggedright
Action
\end{minipage} & \begin{minipage}[b]{\linewidth}\raggedright
Trace
\end{minipage} \\
\midrule\noalign{}
\endhead
\bottomrule\noalign{}
\endlastfoot
10 min & Décrire la personne, son besoin et l'action principale &
Intention en trois phrases \\
15 min & Comparer deux pistes et choisir un parcours court & Choix,
raison et limite \\
20 min & Définir contenu, geste, contribution 3D, HTML et capacité &
Carte du projet et critères \\
15 min & Choisir l'expérience, la remise et les points de revue & Plan
de six heures solo \\
\end{longtable}
}

L'idée appartient à l'élève. La contribution 3D sert un usage expliqué
et le contenu conserve une alternative HTML. Les briques fournies, les
modifications et le raccordement sont attribués séparément.

\section{Six heures solo, suivies de la
clôture}\label{six-heures-solo-suivies-de-la-cluxf4ture}

Distribuer 1 h de préparation, 2 h d'adaptation, 2 h
d'intégration/expérience et 1 h de préparation de remise. Dans le
calendrier, cela correspond aux derniers modules. La séance de clôture
suit ce travail ; aucune heure accompagnée supplémentaire n'est ajoutée.

\section{Une heure de clôture}\label{une-heure-de-cluxf4ture}

{\def\LTcaptype{none} % do not increment counter
\begin{longtable}[]{@{}
  >{\raggedleft\arraybackslash}p{(\linewidth - 4\tabcolsep) * \real{0.4000}}
  >{\raggedright\arraybackslash}p{(\linewidth - 4\tabcolsep) * \real{0.3000}}
  >{\raggedright\arraybackslash}p{(\linewidth - 4\tabcolsep) * \real{0.3000}}@{}}
\toprule\noalign{}
\begin{minipage}[b]{\linewidth}\raggedleft
Durée
\end{minipage} & \begin{minipage}[b]{\linewidth}\raggedright
Action
\end{minipage} & \begin{minipage}[b]{\linewidth}\raggedright
Examen
\end{minipage} \\
\midrule\noalign{}
\endhead
\bottomrule\noalign{}
\endlastfoot
10 min & Montrer besoin et parcours & Usage explicite \\
20 min & Démontrer geste, sélection, retour, statique et HTML &
Comportement observé \\
15 min & Expliquer une modification et rejouer une expérience & Données,
paramètres et limite \\
10 min & Appliquer une notion à un autre contexte & Transfert nouveau \\
5 min & Conserver acquis observés et prochaine question & Bilan
sélectionné \\
\end{longtable}
}

\subsection{Script enseignant --- séquence de 5
minutes}\label{script-enseignant-suxe9quence-de-5-minutes-8}

Les pauses et intonations sont des propositions de lecture.

\textbf{0:00--1:00.} « Je veux partir de ton intention. À qui sert ce
frontend, et quelle action cette personne doit-elle pouvoir accomplir ?
{[}Pause.{]} Montre-moi ce parcours avant de parler des bibliothèques. »

\textbf{1:00--2:00.} « Maintenant, choisis une modification qui
t'appartient. Où est-elle dans le fichier ? Qu'as-tu gardé du modèle, et
quelle aide as-tu reçue ? » Laisser retrouver le code plutôt que réciter
une explication générale.

\textbf{2:00--3:00.} « Prenons ton expérience. {[}Ralentir.{]} Quelle
variable as-tu changée ? Quelles conditions as-tu conservées ?
Montre-moi l'observation et une limite. » Examiner un résultat concret ;
une observation insuffisante conserve son statut.

\textbf{3:00--4:00.} « On va essayer le même contenu avec moins de
mouvement et au clavier. Est-ce que la personne retrouve toujours son
information ? » Laisser réaliser le parcours et consigner un éventuel
défaut, sans l'effacer du bilan.

\textbf{4:00--5:00.} « Tu pourrais réutiliser cette notion ailleurs.
Propose un exemple qui n'est pas cette scène. Quelle donnée ou quelle
règle garderais-tu ? » Laisser reformuler puis définir une prochaine
vérification utile.

\section{Revue de projet}\label{revue-de-projet}

Demander de retrouver une modification réelle dans les exports,
d'expliquer un événement qui change l'état et de montrer la reprise.
Examiner une expérience comparable plutôt qu'une impression seule.
Donner un exemple nouveau de carte, de galerie ou de capteur pour le
transfert.

Conserver une grille : intention, attribution, état, accessibilité,
expérience, capacité bornée, reprise et transfert. Chaque ligne reçoit
le scénario, le constat et sa limite. L'enseignant garde son retour
séparé de la proposition de l'assistant et de l'empreinte des fichiers.
La publication publique reste une action distincte.

\chapter{Conclusion générale --- une production que l'on peut
expliquer}\label{conclusion-guxe9nuxe9rale-une-production-que-lon-peut-expliquer}

Le parcours laisse une notion, une expérience et un fichier par module.
Le projet final relie les productions au besoin choisi. L'enseignant
examine les résultats, les limites et un transfert. L'assistant aide à
avancer ; l'élève peut montrer ce qu'il a compris, ce qu'il conserve et
ce qu'il veut vérifier ensuite.

Les scripts et les critères de ce manuel préparent cet examen. Les
observations d'apprentissage appartiennent aux séances réellement
effectuées.

\chapter{Annexes}\label{annexes}

\section{A --- Remettre et reprendre un
travail}\label{a-remettre-et-reprendre-un-travail}

Avant l'export, vérifier le module, la tentative, le paramètre conservé,
les conditions et la question ouverte. Choisir les fichiers et
observations à remettre. Conserver le notebook et le ZIP dans son espace
personnel.

À l'import, Orbit vérifie les fichiers et les empreintes de l'archive
compatible. Le contenu importé reste déclaré externe. Une nouvelle
version conserve l'ancienne ; son partage est choisi séparément.
Calculer une empreinte permet de comparer le contenu, avec une
explication distincte pour l'exécution et la compréhension.

Si la sauvegarde locale échoue, lire le message durable et exporter la
session avant de fermer. Lors d'une reprise de session, choisir à
nouveau les permissions et les éléments partagés. Un ancien résultat de
l'assistant ne doit pas écraser une révision nouvelle.

Pour Colab, distinguer relecture dans le même runtime, namespace neuf et
redémarrage réel. Retrouver les fichiers sauvegardés après un
redémarrage avant de relancer l'expérience. Un ZIP avec une tentative
réutilisée et un contenu différent demande une correction de la
tentative, pas une importation ambiguë.

\section{B --- Les vingt-cinq outils dans le contexte
pédagogique}\label{b-les-vingt-cinq-outils-dans-le-contexte-puxe9dagogique}

{\def\LTcaptype{none} % do not increment counter
\begin{longtable}[]{@{}
  >{\raggedright\arraybackslash}p{(\linewidth - 2\tabcolsep) * \real{0.5000}}
  >{\raggedright\arraybackslash}p{(\linewidth - 2\tabcolsep) * \real{0.5000}}@{}}
\toprule\noalign{}
\begin{minipage}[b]{\linewidth}\raggedright
Fonction
\end{minipage} & \begin{minipage}[b]{\linewidth}\raggedright
Outils
\end{minipage} \\
\midrule\noalign{}
\endhead
\bottomrule\noalign{}
\endlastfoot
Découverte et méthode & \texttt{orbit\_\allowbreak{}get\_\allowbreak{}capabilities},
\texttt{orbit\_\allowbreak{}get\_\allowbreak{}research\_\allowbreak{}protocol} \\
Context & \texttt{orbit\_\allowbreak{}sanity\_\allowbreak{}initial\_\allowbreak{}context},
\texttt{orbit\_\allowbreak{}sanity\_\allowbreak{}read\_\allowbreak{}entries} \\
Question et mission de recherche &
\texttt{orbit\_\allowbreak{}get\_\allowbreak{}research\_\allowbreak{}request},
\texttt{orbit\_\allowbreak{}get\_\allowbreak{}mission\_\allowbreak{}summary},
\texttt{orbit\_\allowbreak{}list\_\allowbreak{}research\_\allowbreak{}points} \\
Sources & \texttt{orbit\_\allowbreak{}search\_\allowbreak{}sources},
\texttt{orbit\_\allowbreak{}read\_\allowbreak{}source\_\allowbreak{}record} \\
Proposition de recherche & \texttt{orbit\_\allowbreak{}present\_\allowbreak{}research} \\
Preuves et relations & \texttt{orbit\_\allowbreak{}classify\_\allowbreak{}evidence},
\texttt{orbit\_\allowbreak{}compare\_\allowbreak{}claims}, \texttt{orbit\_\allowbreak{}find\_\allowbreak{}relations},
\texttt{orbit\_\allowbreak{}resolve\_\allowbreak{}hold}, \texttt{orbit\_\allowbreak{}trace\_\allowbreak{}impact} \\
Mission et méthode pédagogique & \texttt{orbit\_\allowbreak{}get\_\allowbreak{}learning\_\allowbreak{}mission},
\texttt{orbit\_\allowbreak{}get\_\allowbreak{}learning\_\allowbreak{}protocol} \\
Activité et expérience & \texttt{orbit\_\allowbreak{}plan\_\allowbreak{}learning\_\allowbreak{}activity},
\texttt{orbit\_\allowbreak{}prepare\_\allowbreak{}experiment} \\
Artefact et aide & \texttt{orbit\_\allowbreak{}read\_\allowbreak{}learning\_\allowbreak{}artifact},
\texttt{orbit\_\allowbreak{}get\_\allowbreak{}learning\_\allowbreak{}support} \\
Revue et transfert & \texttt{orbit\_\allowbreak{}check\_\allowbreak{}understanding},
\texttt{orbit\_\allowbreak{}prepare\_\allowbreak{}transfer} \\
Proposition et journal & \texttt{orbit\_\allowbreak{}present\_\allowbreak{}learning\_\allowbreak{}work},
\texttt{orbit\_\allowbreak{}get\_\allowbreak{}learning\_\allowbreak{}journal} \\
\end{longtable}
}

Les commandes humaines restent utilisables lorsque WebMCP est
indisponible. Le registre réellement découvert et la permission courante
gouvernent l'appel. Les références privées exigent la sélection, le
contexte et la révision appropriés. \texttt{PRESENTED} indique un dépôt
pour revue. \texttt{READY} indique une ressource ou un calcul
disponible. Les états \texttt{CONSENT\_\allowbreak{}REQUIRED},
\texttt{STALE\_\allowbreak{}REVISION}, \texttt{NOT\_\allowbreak{}FOUND} et \texttt{UNAVAILABLE}
indiquent la condition à traiter.

Le contexte pédagogique possède un registre par document. Son
propriétaire de montage est distinct du dossier conservé :
\texttt{register\allowbreak{}Formation\allowbreak{}Tools} attend la fermeture de l'ancien
enregistrement avant de reprendre les noms natifs;
\texttt{unregister\allowbreak{}Formation\allowbreak{}Tools} ferme seulement l'enregistrement du
propriétaire concerné. Ainsi, le nettoyage tardif d'une ancienne vue
garde intacte sa vue successeure. Le changement de module, la révocation
et la fermeture invalident les accès privés. Ces responsabilités
décrivent le code partagé; leur preuve navigateur appartient au reçu du
même payload.

Les deux outils Context lisent le corpus du cours par \textbf{GET} sur
\texttt{/\allowbreak{}api/\allowbreak{}v1/\allowbreak{}course-context/\allowbreak{}outline} et
\texttt{/\allowbreak{}api/\allowbreak{}v1/\allowbreak{}course-context/\allowbreak{}entries}, avec les credentials du
navigateur omis. Le transport ne transmet pas le journal personnel. Le
serveur borne ces lectures publiques et garde son activation explicite;
une passerelle indisponible produit \texttt{UNAVAILABLE}. Ce transport
et la vraie session Sanity du Studio sont deux capacités à qualifier
séparément.

Les trois moteurs sont \texttt{baseline}, \texttt{n} et \texttt{p},
seuls ou par paires sur les mêmes données. Une session neuve n'en
exécute aucun implicitement. Leurs résultats restent séparés ; une
moyenne ou un vote ne désigne pas un gagnant. T et F peuvent coexister,
I peut subsister et ces ensembles ne constituent pas des probabilités.
Une date différente demande une relation de version justifiée ; deux
attributs manquants restent inconnus.

\section{C --- Installer le plugin dans son
Studio}\label{c-installer-le-plugin-dans-son-studio}

Le plugin apporte les vues Orbit · Lab et Orbit · Projects à un Studio
déjà détenu par l'élève. Le Studio hôte fournit projet, dataset, session
et droits. Le plugin contient le cœur partagé ; React et Sanity
proviennent de l'hôte. Le guide d'installation maintenu donne l'archive
et les versions examinées.

Les travaux privés commencent en mémoire ou dans le stockage local
autorisé. La publication Sanity est une action humaine : sélectionner
les artefacts publiables, examiner le JSON exact, la destination et la
révision, puis effectuer la commande avec les droits réels. Le journal
privé et les permissions ne sont pas ce contenu publié. Une compilation
du plugin et un parcours authentifié sont des preuves distinctes.

Le second hôte préparé sous \texttt{/\allowbreak{}formation/\allowbreak{}studio/\allowbreak{}} conserve le
projet et le dataset de l'équipe, avec un autre point de montage. La
session et les droits appartiennent toujours à Sanity. Deux routes sur
la même origine gardent le même périmètre d'utilisateur; elles ne créent
pas une séparation de confidentialité. Installer le plugin dans le
Studio personnel de l'élève utilise son propre projet, dataset et
session. Le transport public du cours et cette session authentifiée
gardent des autorisations distinctes.

\section{D --- Carte du code}\label{d-carte-du-code}

La carte détaillée suit le tableau des modules ci-dessous. Les fichiers
de raccordement sont fournis et attribués au cours. Les dossiers
\texttt{src/\allowbreak{}learning/\allowbreak{}module-N/\allowbreak{}} contiennent les exports réellement
sélectionnés. L'assemblage ne remplace pas une brique absente par une
correction cachée.

{\def\LTcaptype{none} % do not increment counter
\begin{longtable}[]{@{}
  >{\raggedright\arraybackslash}p{(\linewidth - 4\tabcolsep) * \real{0.3333}}
  >{\raggedright\arraybackslash}p{(\linewidth - 4\tabcolsep) * \real{0.3333}}
  >{\raggedright\arraybackslash}p{(\linewidth - 4\tabcolsep) * \real{0.3333}}@{}}
\toprule\noalign{}
\begin{minipage}[b]{\linewidth}\raggedright
Module
\end{minipage} & \begin{minipage}[b]{\linewidth}\raggedright
Fichier et symbole
\end{minipage} & \begin{minipage}[b]{\linewidth}\raggedright
Paramètre de l'essai
\end{minipage} \\
\midrule\noalign{}
\endhead
\bottomrule\noalign{}
\endlastfoot
1 &
\href{../../frontend/module-1/interaction-state.js}{interaction-state.js}
· \texttt{create\allowbreak{}Interaction} & keyboardStep : 0.04 → 0.08 \\
2 &
\href{../../frontend/module-2/exploration-card.js}{exploration-card.js}
· \texttt{create\allowbreak{}Cards} & descriptionPrefix : Mon observation → Ma
décision expliquée \\
3 &
\href{../../frontend/module-3/scene-controller.js}{scene-controller.js}
· \texttt{create\allowbreak{}Scene\allowbreak{}Controller} & cameraDistance : 5 → 7 \\
4 & \href{../../frontend/module-4/inertia.js}{inertia.js} ·
\texttt{create\allowbreak{}Inertia} & damping : 0.65 → 1.3 \\
5 & \href{../../frontend/module-5/transitions.js}{transitions.js} ·
\texttt{create\allowbreak{}Transition} & damping : 12 → 5 \\
6 & \href{../../frontend/module-6/particle-layer.js}{particle-layer.js}
· \texttt{create\allowbreak{}Particle\allowbreak{}Layer} & count : 64 → 128, puis restauration à
64 \\
7 &
\href{../../frontend/module-7/interaction-flow.js}{interaction-flow.js}
· \texttt{create\allowbreak{}Interaction\allowbreak{}Flow} & historyLimit : 32 → 4 \\
8 &
\href{../../frontend/module-8/register-capabilities.js}{register-capabilities.js}
· \texttt{register\allowbreak{}Learning\allowbreak{}Capability} & borne de snapshot : 12000 → 6000
caractères JSON \\
Raccordement fourni &
\href{../../assembly/src/student-frontend.js}{student-frontend.js},
\href{../../assembly/merge_modules.py}{merge\_modules.py} & Les huit ZIP
réels, puis compilation et parcours séparés \\
\end{longtable}
}

\section{E --- Glossaire / Glossary}\label{e-glossaire-glossary}

{\def\LTcaptype{none} % do not increment counter
\begin{longtable}[]{@{}
  >{\raggedright\arraybackslash}p{(\linewidth - 2\tabcolsep) * \real{0.5000}}
  >{\raggedright\arraybackslash}p{(\linewidth - 2\tabcolsep) * \real{0.5000}}@{}}
\toprule\noalign{}
\begin{minipage}[b]{\linewidth}\raggedright
Notion
\end{minipage} & \begin{minipage}[b]{\linewidth}\raggedright
Explication liée au cours
\end{minipage} \\
\midrule\noalign{}
\endhead
\bottomrule\noalign{}
\endlastfoot
État & Données conservées pour décrire un geste, une sélection ou une
vue \\
Événement & Signal reçu par le programme, comme une pression, un
déplacement ou un retour navigateur \\
Capture du pointeur & Acheminement des événements vers l'élément qui
tient le geste, avec un identifiant actif \\
Rendu & Présentation calculée à partir de l'état \\
Identifiant stable & Valeur reliant un même élément à sa carte, sa scène
et sa fiche \\
Coordonnée normalisée & Position relative à la zone ; ici entre 0 et
1 \\
Caméra et projection & Point de vue et transformation de la scène vers
une image \\
Maillage & Géométrie et matériau regroupés pour représenter un objet \\
Matériau & Paramètres qui décrivent la réaction visuelle d'une
surface \\
Raycasting & Recherche d'un objet rencontré par un rayon de sélection \\
Buffer & Zone de données réutilisable, par exemple pour positions de
particules \\
Inertie & Conservation d'une vitesse après relâchement dans ce modèle
d'interface \\
Amortissement & Terme réduisant la vitesse dans le modèle fourni \\
Interpolation & Règle reliant une valeur présente à une cible \\
Ressort & Évolution conservant aussi une vitesse et pouvant dépasser la
cible \\
Boucle de rendu & Cadence qui met à jour les données puis demande leur
dessin \\
Historique borné & Liste diagnostique limitée à un nombre d'entrées \\
Nettoyage & Fermeture des événements et libération des ressources
possédées \\
Révision & Version du travail sur laquelle une proposition ou une
permission s'applique \\
Empreinte & Valeur permettant de comparer l'identité d'un contenu \\
Provenance & Origine, référence, version et statut de lecture d'un
élément \\
Portée & Conditions d'application d'une affirmation \\
HOLD & Suspension justifiée indiquant l'information attendue et la
condition de reprise \\
Transfert & Réemploi d'une notion dans une situation nouvelle \\
\end{longtable}
}

\section{F --- Symboles et unités}\label{f-symboles-et-unituxe9s}

{\def\LTcaptype{none} % do not increment counter
\begin{longtable}[]{@{}
  >{\raggedright\arraybackslash}p{(\linewidth - 2\tabcolsep) * \real{0.5000}}
  >{\raggedright\arraybackslash}p{(\linewidth - 2\tabcolsep) * \real{0.5000}}@{}}
\toprule\noalign{}
\begin{minipage}[b]{\linewidth}\raggedright
Symbole ou champ
\end{minipage} & \begin{minipage}[b]{\linewidth}\raggedright
Sens dans les fichiers du cours
\end{minipage} \\
\midrule\noalign{}
\endhead
\bottomrule\noalign{}
\endlastfoot
\texttt{x}, \texttt{y} & Position normalisée, sans unité de longueur
physique \\
\texttt{vx}, \texttt{vy} & Variation de position normalisée par
seconde \\
\texttt{dt} & Durée d'une étape, en secondes ; le modèle borne les
grandes durées \\
\texttt{k} / \texttt{damping} & Coefficient d'amortissement ; dans la
loi exponentielle, \texttt{k\ ×\ dt} est sans dimension \\
\texttt{value}, \texttt{target} & Valeur d'échelle présente et cible
dans la transition \\
\texttt{velocity} & Variation de l'échelle par seconde \\
\texttt{count}, \texttt{age}, \texttt{lifetime} & Population, âge et
durée de vie d'une particule de démonstration \\
T / I / F & Soutiens, raisons d'indétermination et réfutations
indépendantes pour une affirmation \\
\end{longtable}
}

Les nombres servent ici aux mécanismes d'interface. Une graine
répétable, une sphère ou une trajectoire ne donnent pas une validation
physique du modèle.

\section{G --- Index des notions / Concept
Index}\label{g-index-des-notions-concept-index}

\begin{itemize}
\tightlist
\item
  Accessibilité : Modules 1, 2, 3, 5 ; fiches enseignant ; annexes A et
  C.
\item
  Assistant, autorité et attribution : Module 8 ; partie II ; annexes A,
  B et C.
\item
  Caméra, coordonnées et sélection : Module 3 ; carte du code ;
  glossaire.
\item
  État et événements : Modules 1 et 7 ; fiches enseignant 1 et 7.
\item
  Empreintes, preuve et révision : Module 8 ; qualification ; annexes A
  et B.
\item
  Inertie, vitesse et durée : Module 4 ; symboles et unités.
\item
  Interpolation et ressort : Module 5 ; fiche enseignant 5.
\item
  Particules et nettoyage : Modules 3 et 6 ; fiche enseignant 6.
\item
  Projet, remise et reprise : projet personnel ; fiche enseignant projet
  ; annexe A.
\item
  Transfert : critères de chaque module ; clôture du projet.
\end{itemize}

\section{H --- Sources et
références}\label{h-sources-et-ruxe9fuxe9rences}

Les fichiers du cours et les sources primaires cités dans chaque
chapitre sont recensés dans \texttt{source-map.json}. La carte conserve
leurs empreintes au moment de la composition, leur rôle et leur état.
Les sorties générées utilisent ce contenu commun. Les horaires et
missions restent gouvernés par le catalogue ; la qualification reste
gouvernée par le reçu de livraison.

\begin{itemize}
\tightlist
\item
  \href{../../../../../packages/learning/src/catalog.ts}{Catalogue des
  missions}.
\item
  \href{../../FORMATION_CONTRACT.md}{Contrat du cours}.
\item
  \href{../../WEBMCP.md}{Contrat des outils WebMCP}.
\item
  \href{../../STUDIO_INSTALLATION.md}{Installation du plugin Studio}.
\item
  \href{../../assembly/README.md}{Export et assemblage}.
\item
  \href{https://research.google.com/colaboratory/faq.html}{FAQ
  officielle Colab}.
\item
  \href{https://docs.astro.build/en/basics/astro-components/}{Documentation
  Astro}.
\item
  \href{https://threejs.org/manual/\#en/fundamentals}{Manuel Three.js}.
\item
  \href{https://developer.chrome.com/docs/ai/webmcp/imperative-api}{API
  WebMCP Chrome}.
\item
  \href{https://www.sanity.io/docs/ai/sanity-context}{Sanity Context}.
\end{itemize}

\section{Mot de clôture}\label{mot-de-cluxf4ture}

L'élève choisit une production qu'il peut montrer et une notion qu'il
peut expliquer dans un nouvel exemple. L'enseignant conserve ce qui a
été effectivement examiné et prépare la question suivante. Le manuel et
ses documents web soutiennent ce travail concret ; la revue humaine en
garde la responsabilité.

\chapter{Empreintes intégrales de
lecture}\label{empreintes-intuxe9grales-de-lecture}

Les tableaux et paragraphes de cette édition abrègent les identifiants
longs pour garder leur lecture dans la page. Les valeurs intégrales
ci-dessous sont recopiées sans modification depuis le manuscrit
maintenu; les reçus liés précisent leur objet et leur scénario.

Valeurs intégrales, numérotées par ordre de première apparition dans le
manuscrit :

\begin{Shaded}
\begin{Highlighting}[]
\NormalTok{01  28813ecbf712f4fb717aaa3f93444c7c7e14fbd737f2d3e3ac1ed6ca28f72531}
\NormalTok{02  98a0c6594d176b3cb4fe7b921fd93e945e51d96f3fba3019d86037585062ed87}
\NormalTok{03  689494986b4b6a9b02b5a22a8beb81e8de7cb7388a55ad20a122743d3de7966e}
\NormalTok{04  ea0f93eff48e9dd6a07283f46067eafdc317b352e43409efac8e5e61d4992b90}
\NormalTok{05  f2b0dc6897fae7b56e36836e4019eff06e098bf14d35668fa1df84ec6ff9e973}
\NormalTok{06  7e6c739c0883e39aca491e6f6280cc353daf6745017a3bea5ce690746e4d67a5}
\NormalTok{07  4b06da41e603d63acd33bbf75995a75d824d1bc6}
\NormalTok{08  827528418863211e76307d3e3dbd2a21d17d2864bde0e59f4d8295948897f3ce}
\NormalTok{09  96acc0ddfdd3728426980df39d2a1ccda07cc855eb95f950a552de776a5b2aed}
\NormalTok{10  b202088c8f092630676273c7cdf97d89302983866f70f5b8df3305486d7c3335}
\NormalTok{11  8d80562e069f26bb07ec097ad5763676548bf125d7ab06ce5529f8d590bd1654}
\NormalTok{12  141466603afaca979af3e8ee011391c9516de2a723c52809bdca86f0dbb0b1ba}
\NormalTok{13  423e48783d884b71827bef1e4aae380a3e15ca93735ef1a89423cf54e8a0edd8}
\NormalTok{14  e4215188e3f702b8707d3992aaa0f917bbdca00e58118a0f6953d98697412938}
\NormalTok{15  ed94c13524c078f84c44f074a3034d10850f81470a80603638c3404663b4ba06}
\NormalTok{16  d8ff5bd2c54474cca831101f9992d4b92af273370c6f53dee07cd7b639c2d2ce}
\NormalTok{17  3acb942558e5dc6e42104494238a6191af12c274d607f184843611be3166807b}
\NormalTok{18  6b85b3780c50e10cb02a4e75e2ef553b15989f2487cef500ed15b6d572cc0f45}
\NormalTok{19  1426f7985237d66f60736b2b55e96e8b47aa6d9736e161ce25bef40b461b031b}
\end{Highlighting}
\end{Shaded}


\end{document}
